\documentclass[11pt,a4paper]{article}

% Core packages
\usepackage[utf8]{inputenc}
\usepackage[T1]{fontenc}
\usepackage{amsmath,amssymb,amsthm}
\usepackage{mathtools}
\usepackage{longtable,booktabs,array}
\usepackage{geometry}
\geometry{margin=1in}
\usepackage{graphicx}
\usepackage[round,sort&compress]{natbib}
\usepackage{xcolor}
\usepackage{caption}
\usepackage{hyperref}
\usepackage{cleveref}

% TikZ and PGFPlots for figures
\usepackage{tikz}
\usepackage{pgfplots}
\pgfplotsset{compat=1.18}
\usetikzlibrary{patterns, positioning, arrows.meta, shapes.geometric, backgrounds, fit}

% Misc
\usepackage{enumitem}
\usepackage{microtype}

% Custom commands
\newcommand{\R}{\mathbb{R}}
\DeclareMathOperator{\Var}{Var}
\DeclareMathOperator{\Corr}{Corr}
\DeclareMathOperator{\argmax}{arg\,max}
\DeclareMathOperator{\CosSim}{CosSim}

\title{\textbf{The H-Bar Model: Schema-Coherence Suppression as the\\
Origin of Compositional Generalization Failure}\\[6pt]
\large A Pre-Registered Dynamical-Systems Framework}

\author{
  Basyirin Amsyar bin Basri\\
  Independent Researcher\\
  Petaling Jaya, Malaysia\\
  \texttt{basyirin.basri@gmail.com}
}

\date{\today}

% arXiv metadata (uncomment and update when submitting)
% \usepackage{arxiv}
% \arxiv{XXXX.XXXXX}  % Replace with your arXiv ID after submission

\begin{document}

\maketitle

% ============================================================
\subsection*{Abstract}
% ============================================================

Current training pipelines optimise parametric depth $\delta_A(d,t)$
without formally targeting schema coherence $\sigma_A(d,t)$ --- the degree
to which an agent's representations are restructured around deep governing
principles rather than surface-statistical regularities.
\textbf{Existing empirical results on compositional generalisation
benchmarks suggest that agents exhibiting high parametric depth but
low schema coherence tend to fail systematically} on out-of-distribution
recombination tasks that standard in-distribution metrics do not detect.
No existing framework specifies the mechanisms by which $\sigma_A$ forms,
what suppresses it, or how it interacts with attentional allocation
during training.

We introduce the \textbf{H-Bar Model}, a coupled dynamical-systems
framework with three core contributions:
\begin{enumerate}[nosep]
\item \textbf{The $\sigma$-suppression mechanism:} We prove (Proposition~4.1) that standard SGD dynamics create a stable low-$\sigma_A$ equilibrium (the ``$\sigma$-trap''), explaining why high-ID/low-OOD agents arise systematically
\item \textbf{Coupled ODE framework:} A mathematically grounded dynamical system (Eqs.~14, 28, 29) with proven local existence/uniqueness (Prop.~3.1), forward invariance (Prop.~3.2), and bifurcation analysis (Prop.~4.1--4.4)
\item \textbf{Phase-structured curriculum:} A five-phase training arc derived from the ODE's bifurcation structure, with empirically detectable Phase-2 entry signature (Prediction~9)
\end{enumerate}
Extensions covering attentional fidelity, executive control, metacognition,
social cognition, and multimodal transfer are provided in Appendix~B.

The framework demonstrates formal correspondence with five cognitive
faculty gaps identified by \cite{burnell2026measuringprogressagi} and
\textbf{enables generation of executable benchmark families for each}.
A pre-registered experimental protocol (Section~\ref{sec:empirical-validation})
specifies statistical tests for core predictions. Initial experiments on
SCAN/COGS confirm both H-Bar conditions (additive and multiplicative
$\sigma_A$ coupling) significantly outperform standard SGD on
out-of-distribution accuracy ($d = 9.08$ and $7.57$, respectively; $p < 0.0001$), with Phase~2
entry detected early in training. If $\sigma_A$-suppression is the
mechanistic origin of compositional generalisation failure, capable-agent
training requires structural revision beyond scale and curriculum
ordering alone.

\bigskip
\noindent\rule{\linewidth}{0.4pt}

% ============================================================
\section{Introduction}\label{sec:introduction}
% ============================================================

The standard framing of systematic compositionality failure is precise
and reproducible. Agents trained to high in-distribution accuracy on
sequence and language tasks fail near-completely when test distributions
require zero-shot recombination of primitives trained in isolation.
\cite{lake2018generalizationwithoutsystematicity} demonstrated this on
SCAN --- seq2seq models achieve above 99\% accuracy on random test splits
and below 2\% on the add-primitive split. The broader benchmark
literature is consistent: COGS \citep{kim2020cogscompositionalgeneralizationchallenge}
finds 96--99\% in-distribution accuracy alongside 16--35\% compositional
generalisation accuracy; CFQ \citep{keysers2020measuringcompositionalgeneralizationcomprehensive}
shows performance degradation growing monotonically with compound
divergence; PCFG-SET \citep{hupkes2020compositionalitydecomposed} documents
failures on productivity, systematicity, and substitutivity tests across
RNNs, CNNs, and Transformers.

The structural argument follows directly. Current training pipelines have
one formal optimisation target: minimise expected loss over the training
distribution. This objective efficiently increases one variable ---
\textbf{parametric depth} $\delta_A(d,t)$ --- but has no formal
mechanism for increasing a second, independent variable: \textbf{schema
coherence} $\sigma_A(d,t)$, the degree to which representations are
restructured around deep governing principles.

The H-Bar Model V1.0 formalised this asymmetry. V2.0 extended it to
cover four additional cognitive faculties --- Attention, Executive
Functions, Metacognition, and Social Cognition --- each with formal
variables, ODEs, and benchmark generation protocols. V3.0 extended all
variables to a domain $\times$ modality product space and formalised
benchmark validity as a measurable object. V3.0+ adds a reliability
function and three practical protocols that close the remaining
measurement gaps.

\subsection{Central Claim}\label{sec:central-claim}

\begin{quote}
Training pipelines that optimise $\delta_A(d,t)$ without formally
targeting $\sigma_A(d,t)$ are predicted to systematically produce
agents that pass in-distribution evaluation while failing
out-of-distribution recombination --- not because they lack depth,
but because their training regimes suppress the schema crystallisation
that converts parametric depth into principled generalisation capacity.
Furthermore, the suppression mechanism extends to attentional allocation
($\alpha_A$), executive control ($\Xi_A$), metacognitive self-modelling
($\hat M_A$), and inter-agent schema communication ($\mu_{AB}$) ---
all of which interact multiplicatively with $\sigma_A$ through the
formal coupling terms derived below.
\end{quote}

\paragraph{Scope qualification.}
Our central claim applies to agents trained on compositional
generalisation benchmarks (SCAN, COGS, PCFG-SET) using standard
gradient-based optimisation. We do \emph{not} claim that
$\sigma_A$-targeting is the \textbf{only} path to compositional
success. Large-scale pretraining \citep{kaplan2020scaling},
in-context learning \citep{min2022rethinking}, and dual-process
architectures \citep{noviello2025neuroscienceinspiredualprocess} may achieve similar ends via
different mechanisms. The H-Bar framework provides one
\emph{sufficient} mechanism (schema-coherence cultivation) and
predicts its empirical signatures; it does not claim \emph{necessary}
exclusivity.

\subsection{Cognitive Faculty Alignment}\label{sec:cognitive-faculty-alignment}

The framework demonstrates formal correspondence with the five cognitive
faculties identified as evaluation gaps by \cite{burnell2026measuringprogressagi}:

\begin{table}[htbp]
\centering
\caption{Cognitive faculty alignment with H-Bar model variables.}
\label{tab:faculty-alignment}
\begin{tabular}{>{\raggedright}p{3.2cm} >{\raggedright}p{3.5cm} >{\raggedright\arraybackslash}p{7cm}}
\toprule
\textbf{Faculty} & \textbf{Primary Variable} & \textbf{Mechanism} \\
\midrule
\textbf{Learning} & $\sigma_A(d,t)$, $\delta_A(d,t)$ & OOD gap as schema proxy; compositional generalisation \\
\textbf{Metacognition} & $\hat M_A(d,t)$, $\zeta_A(d,t)$ & Self-model accuracy; calibration error dynamics \\
\textbf{Attention} & $\alpha_A(d,t)$, $CA(d,t)$ & Attentional fidelity to generative structure \\
\textbf{Executive Functions} & $\Xi_A^P$, $\Xi_A^I$, $\Xi_A^F$ & Planning, inhibition, cognitive flexibility \\
\textbf{Social Cognition} & $\mu_{AB}$, $\tau_A(B,d,t)$, $\Sigma_{A,B}$ & Schema legibility, theory of mind, collective field \\
\bottomrule
\end{tabular}
\end{table}

\subsection{Argument Roadmap}\label{sec:argument-roadmap}

The central claim (Section~\ref{sec:central-claim}) is established through the
following proof structure, with each step mapped to a paper section:

\textbf{THESIS:} $\delta$-only optimization $\Rightarrow$ $\sigma$-suppression $\Rightarrow$ OOD failure

\textbf{PROOF STRUCTURE:}
\begin{itemize}[nosep]
\item \textbf{Step 1 (Section~\ref{sec:core-framework}):} Define observables ($\delta$, $\beta$, $\sigma$) and show $\sigma$ is latent
  \begin{itemize}[nosep]
  \item Lemma~3.1: $\sigma$ cannot be inferred from $\delta$ alone [Proposition~3.6--3.7]
  \end{itemize}
\item \textbf{Step 2 (Section~\ref{sec:cognitive-dimensions}):} Show standard SGD dynamics suppress $\sigma$
  \begin{itemize}[nosep]
  \item Theorem~4.1: $\partial\sigma/\partial t|_{\text{SGD}} < 0$ when $\Omega_{AI}$ is high [Eq.~\eqref{eq:sigma-ode}]
  \end{itemize}
\item \textbf{Step 3 (Section~\ref{sec:phase-structure}):} Characterize failure mode phase structure
  \begin{itemize}[nosep]
  \item Proposition~4.1--4.4: Phase~2 entry is necessary and sufficient for OOD success
  \end{itemize}
\item \textbf{Step 4 (Section~\ref{sec:predictions}):} Empirical signatures detectable in benchmarks
  \begin{itemize}[nosep]
  \item Predictions~1--9: Each maps to one aspect of thesis [Section~\ref{sec:statistical-framework}]
  \end{itemize}
\end{itemize}

This roadmap is referenced at the start of each major section to maintain
argumentative coherence. The benchmark specification and target validity
components are provided in Section~\ref{sec:empirical}.

\noindent\rule{\linewidth}{0.4pt}

% ============================================================
\section{Related Work and Gap Analysis}\label{sec:related-work}
% ============================================================

\subsection{Curriculum Learning}\label{sec:curriculum-learning}

\cite{bengio2009curriculumlearning} established the foundational result:
easy-to-hard sample ordering improves generalisation by providing
representational scaffolding. Subsequent work elaborated difficulty
measurement \citep{kumar2010selfpacedlearning,portelas2020automaticcurriculumlearningdeep,
narvekar2020curriculumlearningreinforcementlearning} and pacing schedules.
Every variant maximises the rate at which $\delta_A$ increases.

\textbf{Gap statement.} Curriculum learning provides training-sequence
optimisation for depth growth rate but has no formal account of
$\sigma_A(d,t)$ dynamics, no triggering conditions for schema
crystallisation, and no prescription for the qualitative shift in
training regime that $\sigma_{\text{critical}}$-crossing requires. The
H-Bar phase structure formalises all three.

\subsection{Compositional Generalisation}\label{sec:compositional-generalisation}

The SCAN/COGS/CFQ/PCFG-SET benchmark family collectively documents the
high-$\delta_A$/low-$\sigma_A$ failure mode empirically. The consistent
diagnosis: models encode statistical regularities rather than the
compositional rules governing the distribution.

\textbf{Gap statement.} The compositional generalisation literature
characterises the failure mode and provides measurement proxies for
$\sigma_{\text{critical}}$-crossing, but does not formalise $\sigma_A$
as the training variable whose dynamics the failure reveals, does not
identify what suppresses it, and does not specify how to design training
regimes that reliably induce schema crystallisation.

\paragraph{Recent Work Differentiation}

\textbf{\cite{noviello2025neuroscienceinspiredualprocess} --- MIRAGE dual-process model.}
Noviello et al.\ argue that compositional generalisation requires a
separate symbolic System-2 module. The H-Bar framework demonstrates that
this structural requirement is unnecessary by formalising compositional
capacity as a developmental trajectory within a single ODE system.
Specifically, schema coherence $\sigma_A(d,t)$ grows continuously
through the $\sigma_A$ ODE (Eq.~\ref{eq:sigma-ode}), where the
$\alpha_A$ gating term directs neural training effort toward
compositional structure without requiring a separate symbolic substrate.
The $\sigma_{\text{critical}}$ bifurcation (Eq.~\ref{eq:sigma-critical})
defines a threshold crossing that is achievable entirely within the
neural parameter space --- no symbolic overlay is invoked.

\textbf{\cite{li2023slogstructuralgeneralizationbenchmark} --- SLOG benchmark recursion failures.}
Li et al.\ document persistent failures on unseen recursion depths.
H-Bar predicts these failures from the $\sigma_A$ ODE as a Phase~1
diagnostic: agents with $\sigma_A \approx 0$ cannot transfer recursive
patterns to unseen depths because the schema growth term is gated by
$\alpha_A$, which is suppressed when surface-reward pressure
$R_A^{\text{surface}}$ is high.

\textbf{\cite{bruns2025exploringcompositionalgeneralizationin} --- RASP structural representability proof.}
Bruns proves via RASP that transformers possess the structural capacity
to represent compositional operations. H-Bar's framework is
complementary: the RASP proof establishes that
$\sigma_{\text{critical}}$-crossing is theoretically possible, while
H-Bar's ODE system explains why standard training does not achieve it.

\textbf{\cite{jiang2022mutualexclusivity} --- Mutual exclusivity training on COGS.}
Jiang et al.\ demonstrate that mutual exclusivity (ME) training pushes
COGS lexical scores substantially higher while structural
compositionality splits remain below 1\%. H-Bar's $\sigma_A$ ODE
predicts this dissociation precisely: ME training increases
$\delta_A$ at the lexical level but does not gate through $\alpha_A$
for structural compositional regularities, so
$\rho \cdot P_A \cdot \alpha_A \cdot (1 - \sigma_A) \approx 0$
regardless of $P_A$ magnitude.

\textbf{\cite{swayamdipta2020datasetcartography} --- Dataset cartography curricula.}
Swayamdipta et al.\ show that dataset cartography yields performance
gains without closing structural compositionality gaps. This result is
predicted by H-Bar's gap statement: curriculum learning optimises
$\eta(d,t)$ but has no formal mechanism for increasing $\sigma_A$.

\subsection{Continual Learning and Decay}\label{sec:continual-learning}

\cite{mccloskey1989catastrophicinterference},
\cite{goodfellow2013empiricalinvestigationcatastrophicforgetting}, and
\cite{kirkpatrick2017overcomingcatastrophicforgetting} address parametric
overwriting under sequential training. Every mitigation strategy targets
preservation or recovery of parametric content.

\textbf{Gap statement.} The continual learning literature formally
accounts for parametric decay $\lambda_c$ under sequential training but
has no formal object for the domain frontier $\Delta(d,t)$ and no
mechanism for representing frontier obsolescence $\lambda_f(d,t)$ as a
distinct decay process with different intervention implications.

\textbf{Connection to schema coherence.} Parametric decay $\lambda_c$ in
continual learning is the temporal analogue of schema coherence suppression
$\sigma_A$ in compositional generalisation. Just as $\lambda_c$ causes
forgetting of previous tasks, $\sigma_A$-suppression causes ``forgetting''
of compositional structure. The H-Bar framework unifies these through the
schema-mediated decay reduction mechanism (Equation~\ref{eq:effective-decay}):
schema coherence acts as a protective factor against both forms of decay,
with $\gamma_\sigma$ modulating the degree of protection.

\subsection{Causal and Structured Representation Learning}\label{sec:causal-repr}

\cite{scholkopf2021towardcausalrepresentationlearning} establish that
causal representations support OOD generalisation.
\cite{mohan2024structureindrlsurveyopenproblems} document seven structural
prior incorporation patterns. \cite{torresan2024disentangledrepresentationscausalcognition}
distinguish weak from strong disentanglement.

\textbf{Gap statement.} Causal and structured representation learning
specify the target state corresponding to high $\sigma_A(d,t)$ but do
not formalise $\sigma_A$ as a dynamic scalar with its own ODE, do not
model the training process that builds it or the AI bypass risk
$\Omega_{AI}(d,t)$ that suppresses it, and do not connect it to
attentional allocation, executive control, metacognitive self-modelling,
or collective schema communication.

\subsection{Cognitive Evaluation of AI Systems}\label{sec:cognitive-eval}

\cite{burnell2026measuringprogressagi} propose a Cognitive Taxonomy of
ten faculties and identify Learning, Metacognition, Attention, Executive
Functions, and Social Cognition as having large evaluation coverage gaps.
\cite{chollet2019measureintelligence} proposes ARC as a generalisation
benchmark. \cite{morris2023levelsagi} establish a Levels of AGI framework.

\textbf{Gap statement.} Existing cognitive evaluation frameworks specify
what should be measured but provide no formal theoretical grounding for
why specific task designs isolate specific faculties. The H-Bar Benchmark
Protocol (Section~\ref{sec:benchmark-protocol}) provides that grounding
by deriving benchmark designs directly from formal variable structures.

\subsection{Synthesis --- The Five-Gap Map}\label{sec:five-gap-map}

\begin{center}
\begin{tabular}{>{\raggedright}p{3.5cm} >{\raggedright}p{3.5cm} >{\raggedright\arraybackslash}p{6.5cm}}
\toprule
\textbf{Literature} & \textbf{H-Bar Variable Addressed} & \textbf{H-Bar Variable Missing} \\
\midrule
Curriculum Learning & $\delta_A$ growth rate & $\sigma_A$ dynamics, $\alpha_A$, $\Xi_A$ \\
Compositional Generalisation & $\sigma_A$ failure mode (empirical) & $\sigma_A$ formation mechanism, suppression \\
Continual Learning & $\lambda_c$ (parametric decay) & $\lambda_f$ (frontier obsolescence), $\sigma_A$ coupling \\
Causal/Structured Repr. & $\sigma_A$ target state & $\sigma_A$ developmental trajectory, $\hat M_A$, $\mu_{AB}$ \\
Cognitive Evaluation & Faculty identification & Formal theoretical grounding for benchmark design \\
\bottomrule
\end{tabular}
\end{center}

\paragraph{Comparison with Contemporary Frameworks.}
Table~\ref{tab:framework-comparison} positions H-Bar relative to recent
approaches addressing compositional generalisation failure.

\begin{table}[htbp]
\centering
\caption{\textbf{Comparison of compositional generalisation frameworks.}}
\label{tab:framework-comparison}
\small
\begin{tabular}{>{\raggedright\arraybackslash}p{2.0cm}>{\raggedright\arraybackslash}p{2.5cm}>{\raggedright\arraybackslash}p{2.5cm}>{\raggedright\arraybackslash}p{5.0cm}}
\toprule
\textbf{Framework} & \textbf{Failure Mode} & \textbf{Mechanism} & \textbf{Relation to H-Bar} \\
\midrule
MIRAGE \citep{noviello2025neuroscienceinspiredualprocess}
  & System-1/System-2 competition
  & Dual-process architecture
  & Complementary: H-Bar formalises the \emph{dynamics} of System-2 ($\sigma_A$) growth \\

RASP \citep{bruns2025exploringcompositionalgeneralizationin}
  & Capacity gap
  & Restricted access sequences
  & H-Bar explains \emph{when} capacity is utilised (high $\sigma_A$) vs.\ wasted (low $\sigma_A$) \\

Mutual Exclusivity \citep{jiang2022mutualexclusivity}
  & Lexical overlap interference
  & Augmentation-induced sparsity
  & Special case: ME targets $\sigma_A$-indirectly via data; H-Bar targets $\sigma_A$-directly via training dynamics \\

Dataset Cartography \citep{swayamdipta2020datasetcartography}
  & Example difficulty mismatch
  & Curriculum ordering
  & H-Bar extends: curriculum ordered by \emph{phase} ($\sigma_A$-dependent), not just difficulty \\

\textbf{H-Bar (Ours)}
  & \textbf{$\sigma_A$-suppression trap}
  & \textbf{Coupled ODE + phase curriculum}
  & \textbf{Unifies: predicts when prior methods succeed/fail} \\
\bottomrule
\end{tabular}
\end{table}

\paragraph{Distinctive contribution.}
Unlike frameworks that propose \emph{architectural} solutions (MIRAGE,
RASP) or \emph{data-centric} solutions (Mutual Exclusivity, Dataset
Cartography), H-Bar proposes a \emph{dynamical-systems} explanation:
compositional failure is not merely a capacity or data problem, but a
\textbf{bifurcation phenomenon}. Standard training creates a stable
low-$\sigma_A$ equilibrium (Proposition~4.2); escape requires crossing
a critical threshold via phase-aware curriculum. This predicts that
architectural and data-centric methods will succeed only when they
incidentally raise $\sigma_A$ above $\sigma_{\text{critical}}$---a
testable meta-prediction (Meta-Prediction~M\ref{sec:meta-M1}).

\noindent\rule{\linewidth}{0.4pt}

% ============================================================
\section{The H-Bar Model --- Core Framework}\label{sec:core-framework}
% ============================================================

The H-Bar Model formalises agent knowledge development as a coupled
dynamical system. This section presents the foundational architecture
--- the three core state variables (depth, breadth, schema coherence),
their ODEs, intersection activation, and the delegation gradient.
Subsequent sections extend this core with additional cognitive
dimensions (Section~\ref{sec:cognitive-dimensions}), multimodal coverage
and benchmark validity (Section~\ref{sec:multimodal}), and reliability
protocols (Section~\ref{sec:reliability}), all of which couple back into
the system established here.

\subsection{The Three Core Knowledge Dimensions}\label{sec:three-core}

Let every domain of knowledge be indexed by $d \in D$, where $D$ is the
set of all knowledge domains. For a given agent $A$ at time $t$, three
orthogonal quantities are tracked per domain.

\subsubsection{Depth \texorpdfstring{$\delta_A(d,t)$}{delta\_A(d,t)}}\label{sec:depth}

The structural complexity and principled coherence of the agent's
internal representation of domain $d$. Not raw parameter volume ---
the degree to which knowledge is organised around deep causal principles
rather than surface-statistical features. Bounded above by the domain
frontier $\Delta(d,t)$:
\begin{equation}\label{eq:depth-bounds}
0 \leq \delta_A(d,t) \leq \Delta(d,t)
\end{equation}

\textbf{Relative depth:}
\begin{equation}\label{eq:relative-depth}
\delta_A^{\text{relative}}(d,t) = \frac{\delta_A(d,t)}{\Delta(d,t)} \in [0,1]
\end{equation}

\subsubsection{Breadth \texorpdfstring{$\beta_A(d,t)$}{beta\_A(d,t)}}\label{sec:breadth}

The agent's functional competence in domain $d$ --- sufficient to
engage primary-domain artefacts --- but without deep principled
organisation. Qualitatively different from depth: more rapidly acquired,
more rapidly lost, and increasingly AI-augmentable.

\subsubsection{Schema Coherence \texorpdfstring{$\sigma_A(d,t)$}{sigma\_A(d,t)}}\label{sec:schema-coherence}

\textbf{Conceptual definition.} Schema coherence $\sigma_A(d,t) \in [0,1]$
quantifies the degree to which an agent's representation of domain $d$
has been restructured around deep governing principles rather than
surface-statistical features. High $\sigma_A$ indicates that the agent's
internal model encodes compositional operators, causal structure, and
systematic rules applicable beyond the training distribution. Low
$\sigma_A$ indicates memorisation of surface patterns that fail to
transfer.

\textbf{Distinguishing $\sigma_A$ from task accuracy.} While $\sigma_A$
correlates with out-of-distribution (OOD) accuracy, it is conceptually
distinct:
\begin{itemize}[nosep]
\item $\sigma_A$ is a property of the \textit{representation} (internal state)
\item OOD accuracy is a property of \textit{behavioural output}
\item Two agents can achieve identical OOD accuracy via different mechanisms (memorisation vs.\ genuine generalisation); only the latter exhibits high $\sigma_A$
\item $\sigma_A$ can be intermediate (partially coherent) while OOD accuracy is binary (correct/incorrect on specific examples)
\end{itemize}

\textbf{Operationalisation challenge.} $\sigma_A$ is not directly
observable from activations alone. We operationalise it through a
two-stage proxy architecture (Section~\ref{sec:online-estimation}) that
combines training-time structural signals with evaluation-time
calibration. This separation ensures that $\hat{\sigma}_A$ captures
representational properties rather than merely recapitulating OOD labels.

\textbf{Unique properties distinguishing $\sigma_A$ from adjacent
constructs:} $\sigma_A$ uniquely (PC1)~encodes compositional
recombination capacity, (PC2)~is normalised against the domain frontier
$\Delta(d,t)$, and (PC3)~has an evaluative function detecting AI bypass
--- properties not shared by Structured Representations, Disentangled
Representations, Causal Representations, or Cognitive Schemas.

\subsubsection{Online Estimation Protocol}\label{sec:online-estimation}

The proxy identification in Equation~\eqref{eq:sgg-proxy} requires
post-evaluation benchmark data, creating a circular dependency: $\sigma_A$
is needed during training to guide the ODE dynamics, but can only be
measured after training completes. We resolve this through a \textbf{two-stage
estimation architecture} that separates training-time operative estimation
from evaluation-time ground-truth calibration.

\textbf{Stage 1 --- Training-time operative estimate ($\tilde{\sigma}_A$).}
During training, $\sigma_A$ is estimated via a \textbf{multi-signal fusion}
combining three training-time observables:

\textit{(i) Gradient--composition alignment (GCA):}
\begin{equation}\label{eq:gca}
g_A(d,t) = \Corr\!\left(\nabla_{\theta}\mathcal{L}_{\text{comp-batch}},\;\nabla_{\theta}\mathcal{L}_{\text{train}}\right) \in [-1,1]
\end{equation}
where $\mathcal{L}_{\text{comp-batch}}$ is loss on a compositional-recombination
mini-batch constructed programmatically from the domain's generative grammar
at training time.

\textit{(ii) Representational-geometry alignment (RGA):}
\begin{equation}\label{eq:rga}
r_A(d,t) = \Corr\!\Big(\text{RDM}_{\text{rep}}(d,t),\;\text{RDM}_{\text{struct}}(d)\Big) \in [-1,1]
\end{equation}
where $\text{RDM}_{\text{rep}}$ is the representational dissimilarity matrix
from internal activations and $\text{RDM}_{\text{struct}}$ is the structural
dissimilarity matrix from the causal generative graph.

\textit{(iii) Augmentation consistency (AC):}
\begin{equation}\label{eq:ac}
c_A(d,t) = \frac{1}{|A|}\sum_{a \in A} \CosSim\!\Big(R_{\theta}(x),\;R_{\theta}(a(x))\Big) \in [0,1]
\end{equation}
where $A$ is the set of structure-preserving augmentations (primitive
substitution, argument permutation).

\textit{Fusion:}
\begin{equation}\label{eq:fusion}
\tilde{\sigma}_A(d,t) = w_g \cdot \max(0, g_A) + w_r \cdot \max(0, r_A) + w_c \cdot c_A
\end{equation}
with default weights $w_g = 0.4$, $w_r = 0.35$, $w_c = 0.25$ (tunable per
domain). All three signals are computable at any training checkpoint from
forward passes and gradient statistics --- no external benchmarks required.

\textbf{Stage 2 --- Evaluation-time ground-truth calibration ($\hat{\sigma}_A$).}
\begin{equation}
\hat{\sigma}_A(d,t) = \frac{\text{Acc}_{\text{OOD}}(d,t)}{\text{Acc}_{\text{In}}(d,t)}
\label{eq:sgg-proxy}
\end{equation}
computed at evaluation checkpoints on SCAN/COGS/PCFG-SET OOD splits.
Stage 2 serves as the ground-truth calibration for Stage 1. At each
evaluation checkpoint, the Stage 1 estimate $\tilde{\sigma}_A$ is compared
against the Stage 2 value $\hat{\sigma}_A$; if
$|\tilde{\sigma}_A - \hat{\sigma}_A| > 0.15$, the Stage 1 weights
$(w_g, w_r, w_c)$ are updated via ridge regression to minimise the
discrepancy on the evaluation checkpoint history.

\textbf{Operational convention.} Throughout the ODE system
(Equations~\eqref{eq:depth-decay}, \eqref{eq:sigma-ode},
\eqref{eq:discovery-rate}, and all coupled equations), the operative
value of $\sigma_A$ is $\tilde{\sigma}_A(d,t)$ (Stage 1). The validation
value $\hat{\sigma}_A(d,t)$ (Stage 2) is used at evaluation checkpoints
for weight recalibration and for reporting. This convention ensures that
$\sigma_A$ functions as a genuine operative state variable evolving
continuously during training, not a post hoc summary statistic.

\textbf{Proposition 3.6 (Estimator Consistency).}
\textit{Let $\delta_\sigma(d,t) = \tilde{\sigma}_A(d,t) - \hat{\sigma}_A(d,t)$
denote the Stage 1--Stage 2 discrepancy. Under the assumption that the
generative grammar $G(d)$ correctly models the domain's compositional
structure, the multi-signal fusion estimator satisfies:
\begin{equation}\label{eq:estimator-bound}
\mathbb{E}[|\delta_\sigma|] \leq \frac{C_1}{\sqrt{n_{\text{batch}}}} + C_2 \cdot \epsilon_{\text{model}}
\end{equation}
where $n_{\text{batch}}$ is the number of compositional probe samples,
$\epsilon_{\text{model}}$ is the model misspecification error, and $C_1, C_2$
are constants depending on the domain parameters.}

\textit{Proof sketch.} The bound follows from standard concentration
inequalities for U-statistics (for the correlation-based signals) and
Hoeffding's inequality (for the augmentation consistency). The key
requirement is that the compositional probe batch is drawn i.i.d. from
the recombination distribution induced by $G(d)$. $\square$

\textbf{Proposition 3.7 (Error Propagation in ODE System).}
\textit{Let $\delta_\sigma(d,t)$ be the estimation error. The effective
error propagated into the depth ODE \eqref{eq:depth-ode} and schema ODE
\eqref{eq:sigma-ode} is bounded by:
\begin{equation}\label{eq:error-propagation}
|\dot{\delta}_A^{\text{error}}| \leq \lambda_c \cdot \gamma_\sigma \cdot |\delta_\sigma| \cdot \delta_A + \rho \cdot P_A \cdot \alpha_A \cdot |\delta_\sigma|
\end{equation}
If $|\delta_\sigma| < 0.10$, the propagated error is negligible relative
to typical ODE step sizes ($\eta_{\max} \sim O(1)$).}

% --- Figure 1: The H-Bar Conceptual Framework ---
\begin{figure}[htbp]
\centering
\resizebox{0.85\textwidth}{!}{\input{figures/figure1-framework}}
\caption{\textbf{The H-Bar Conceptual Framework.} Knowledge is
structured into vertical parametric depth ($\delta_A$), horizontal
functional breadth ($\beta_A$), and schema coherence ($\sigma_A$) which
acts as the structural integrity required to support cross-domain
intersection activation ($\Psi_A$).}
\label{fig:hbar-framework}
\end{figure}

\subsection{Mastery Set and Breadth Set}\label{sec:mastery-set}

\begin{equation}\label{eq:mastery-set}
M_A(t) = \{d \in D : \delta_A(d,t) > \theta_{\delta} \cdot \Delta(d,t) \text{ AND } \sigma_A(d,t) > \theta_{\sigma}\}
\end{equation}
\begin{equation}\label{eq:breadth-set}
B_A(t) = \{d \in D \setminus M_A(t) : \beta_A(d,t) > \theta_{\beta}\}
\end{equation}

\subsection{Decay Architecture}\label{sec:decay}

\subsubsection{Schema-Mediated Decay Reduction}

Effective decay rate modulated by schema coherence:
\begin{equation}\label{eq:effective-decay}
\lambda_c^{\text{eff}}(d,t) = \lambda_c \cdot (1 - \gamma_\sigma \cdot \sigma_A(d,t))
\end{equation}

\textbf{Design rationale.} The multiplicative form is chosen over alternatives
for three reasons: (a) an additive form $\lambda_c^{\text{eff}} = \lambda_c - \gamma_\sigma \sigma_A$
permits negative decay rates when $\sigma_A > \lambda_c / \gamma_\sigma$, which is
unphysical; (b) the multiplicative form guarantees $\lambda_c^{\text{eff}} \in [0, \lambda_c]$
for all $\sigma_A \in [0,1]$ and $\gamma_\sigma \in [0,1]$; (c) it reduces to the
baseline $\lambda_c$ when $\sigma_A \to 0$, maintaining dimensional consistency.

\subsubsection{Combined Depth Decay}
\begin{equation}\label{eq:depth-decay}
\dot{\delta}_A(d,t)\big|_{\text{decay}} = -\lambda_c \cdot (1 - \gamma_\sigma \cdot \sigma_A(d,t)) \cdot \delta_A(d,t) \cdot (1 - r_A(d,t))
\end{equation}

\subsection{Growth Dynamics}\label{sec:growth}

\subsubsection{Depth Growth ODE}
\begin{equation}\label{eq:depth-ode}
\dot{\delta}_A(d,t) = f_{\text{learn}}(d,t) \cdot \eta(d,t) \cdot T_A(d,t) - \lambda_c \cdot (1 - \gamma_\sigma \cdot \sigma_A(d,t)) \cdot \delta_A(d,t) \cdot (1 - r_A(d,t))
\end{equation}

\textbf{Learning efficiency} --- Gompertz form:
\begin{equation}\label{eq:gompertz}
\eta(d,t) = \eta_{\text{max}} \cdot \exp(-a \cdot \exp(-b \cdot \delta_A^{\text{relative}}(d,t)))
\end{equation}

\textbf{Design rationale.} The Gompertz form is chosen over alternatives
for three reasons: (a) it captures the empirically observed pattern of
rapid initial learning followed by asymptotic approach to the frontier;
(b) unlike logistic growth, it is asymmetric, with the inflection point
at $\delta_A^{\text{relative}} \approx 0.37$ rather than $0.5$, matching
observed learning curves in compositional generalisation tasks; (c) the
double-exponential form ensures $\eta(d,t) > 0$ for all finite depth,
preventing complete learning arrest before frontier attainment.

\subsection{Intersection Activation \texorpdfstring{$\Psi_A(d_1,d_2,t)$}{Psi\_A(d\_1,d\_2,t)}}\label{sec:intersection}

\begin{equation}\label{eq:activation-condition}
I(d_1,d_2)_{\text{active}} \iff \delta_A(d_1,t) > \theta_I \text{ AND } \delta_A(d_2,t) > \theta_I
\end{equation}

\textbf{Effective mastery quality:}
\begin{equation}\label{eq:mastery-quality}
q_A(d,t) = \sigma_A(d,t) \cdot \delta_A^{\text{relative}}(d,t) \in [0,1]
\end{equation}

\textbf{Discovery rate:}
\begin{equation}\label{eq:discovery-rate}
\Psi_A(d_1,d_2,t) = \Psi_0 \cdot \phi(d_1,d_2) \cdot \sqrt{q_A(d_1,t) \cdot q_A(d_2,t)}
\end{equation}

The multiplicative $\sigma_A \cdot \sigma_A$ dependence is the
mechanism's theoretical core: an agent with high $\delta_A$ but low
$\sigma_A$ in one mastery domain shows disproportionately lower $\Psi_A$
than an additive model predicts. This is \textbf{Prediction~6}
(Section~\ref{sec:predictions}).

\subsection{Delegation Gradient \texorpdfstring{$D^*(d,t)$}{D*(d,t)}}\label{sec:delegation}

\begin{equation}\label{eq:delegation-hbar}
D^*_{\text{H-Bar}}(d,t) = \{s \in d : \delta_{AI}(s,t) \geq \delta_A(s,t) \text{ AND } \sigma_A(d,t) \geq \sigma_{\text{critical}}\}
\end{equation}

\noindent\rule{\linewidth}{0.4pt}

% ============================================================
\subsection{Mathematical Foundations}\label{sec:math-foundations}
% ============================================================

The coupled dynamical system defined by Equations~\eqref{eq:depth-ode},
\eqref{eq:sigma-ode}, \eqref{eq:alpha-ode}, \eqref{eq:executive-ode},
\eqref{eq:self-model-ode}, and their multimodal extensions forms a
high-dimensional nonlinear ODE system. This subsection establishes the
mathematical well-posedness of the system.

\subsubsection{State Space and Vector Field}\label{sec:state-space}

Let the state vector for agent $A$ across $|D|$ domains and $|M|$
modalities be:
\begin{equation}\label{eq:state-vector}
\mathbf{x}_A(t) = \left(\{\delta_A(d,m,t)\}, \{\sigma_A(d,m,t)\}, \{\alpha_A(d,m,t)\}, \{\hat{M}_A(d,m,t)\}, \Xi_A^P(t), \Xi_A^I(t), \Xi_A^F(t)\right)
\end{equation}
with dimension $N = |D|\times|M|\times 4 + 3$. The natural state space is:
\begin{equation}\label{eq:state-space}
\mathcal{X} = [0, \Delta_{\max}]^{ |D||M| } \times [0,1]^{3|D||M| + 3}
\end{equation}

\textbf{Proposition 3.1 (Local Existence and Uniqueness).}
\textit{For any initial condition $\mathbf{x}_A(0) \in \text{int}(\mathcal{X})$,
the coupled ODE system admits a unique solution $\mathbf{x}_A(t)$ defined
on a maximal interval $[0, T_{\max})$ where $T_{\max} > 0$.}

\textit{Proof sketch.} The vector field $\mathbf{f}(\mathbf{x}_A)$ defined
by the right-hand sides of the ODE system is composed of: (a) polynomial
terms in state variables (e.g., $\sigma_A \cdot \delta_A$, $\alpha_A \cdot (1-\sigma_A)$),
(b) exponential terms (Gompertz efficiency $\eta(d,t)$), and (c) bounded
rational functions (e.g., $(1 - R_0^{-1})$). All component functions are
$C^\infty$ on $\text{int}(\mathcal{X})$, hence locally Lipschitz. By the
Picard--Lindel\"of theorem, a unique local solution exists. $\square$

\textbf{Proposition 3.2 (Forward Invariance and Boundedness).}
\textit{The state space $\mathcal{X}$ is forward-invariant under the dynamics:
if $\mathbf{x}_A(0) \in \mathcal{X}$, then $\mathbf{x}_A(t) \in \mathcal{X}$
for all $t \in [0, T_{\max})$. Furthermore, if $\Delta(d,t)$ is uniformly
bounded by $\Delta_{\max} < \infty$, then $T_{\max} = \infty$ (global existence).}

\textit{Proof sketch.} We verify boundary behavior for each variable:
\begin{itemize}[nosep]
\item $\delta_A = 0$: $\dot{\delta}_A \geq 0$ (growth term dominates when depth is zero)
\item $\delta_A = \Delta(d,t)$: $\dot{\delta}_A \leq 0$ (approaching frontier reduces growth)
\item $\sigma_A = 0$: $\dot{\sigma}_A \geq 0$ (growth term $\rho \cdot P_A \cdot \alpha_A > 0$)
\item $\sigma_A = 1$: $\dot{\sigma}_A \leq 0$ (the $(1-\sigma_A)$ factor vanishes)
\item $\alpha_A = 0$: $\dot{\alpha}_A \geq 0$ (growth term $\gamma \cdot C_A > 0$)
\item $\alpha_A = 1$: $\dot{\alpha}_A \leq 0$ (the $(1-\alpha_A)$ factor vanishes)
\item $\hat{M}_A = 0$: $\dot{\hat{M}}_A \geq 0$ (driven toward $\sigma_A \geq 0$)
\item $\hat{M}_A = 1$: $\dot{\hat{M}}_A \leq 0$ (mean-reverting toward $\sigma_A \leq 1$)
\item $\Xi_A^{\cdot} = 0$: $\dot{\Xi}_A^{\cdot} \geq 0$ (driven toward targets $P^*, I^*, F^* \geq 0$)
\item $\Xi_A^{\cdot} = 1$: $\dot{\Xi}_A^{\cdot} \leq 0$ (mean-reverting)
\end{itemize}
All boundaries are either reflecting or the vector field points inward.
Since $\mathcal{X}$ is compact when $\Delta_{\max} < \infty$, solutions
cannot escape to infinity, implying $T_{\max} = \infty$. $\square$

\subsubsection{Timescale Separation}\label{sec:timescale}

The system contains rate constants spanning multiple orders of magnitude:
$\eta_{\max} \sim O(1)$ (depth growth), $\rho \sim O(10^{-1})$ (schema growth),
$\gamma \sim O(10^{-1})$ (attention), $\nu_M \sim O(10^{-2})$ (self-model),
$\kappa_P, \kappa_I, \kappa_F \sim O(10^{-2})$ (executive control).

\textbf{Proposition 3.3 (Fast-Slow Decomposition).}
\textit{Define the timescale parameter
$\epsilon = \max(\nu_M, \kappa_P, \kappa_I, \kappa_F) / \min(\eta_{\max}, \rho, \gamma) \ll 1$.
Decompose $\mathbf{x}_A = (\mathbf{x}_A^{\text{fast}}, \mathbf{x}_A^{\text{slow}})$ where
$\mathbf{x}_A^{\text{fast}} = (\delta_A, \sigma_A, \alpha_A)$ and
$\mathbf{x}_A^{\text{slow}} = (\hat{M}_A, \Xi_A^P, \Xi_A^I, \Xi_A^F)$.
Then the system admits a slow manifold
$\mathcal{M}_\epsilon = \{ \mathbf{x}_A : \mathbf{x}_A^{\text{fast}} = \mathbf{h}(\mathbf{x}_A^{\text{slow}}, \epsilon) \}$
that is $O(\epsilon)$-close to the critical manifold
$\mathcal{M}_0 = \{ \mathbf{x}_A : \dot{\mathbf{x}}_A^{\text{fast}} = 0 \}$,
and the reduced dynamics on $\mathcal{M}_\epsilon$ are given by:
\begin{equation}\label{eq:reduced-dynamics}
\dot{\mathbf{x}}_A^{\text{slow}} = \epsilon \cdot \mathbf{g}(\mathbf{h}(\mathbf{x}_A^{\text{slow}}, \epsilon), \mathbf{x}_A^{\text{slow}})
\end{equation}
where $\mathbf{g}$ is the slow subsystem vector field.}

\textit{Proof sketch.} This follows from Fenichel's geometric singular
perturbation theory \citep{fenichel1979geometric}. The critical manifold
$\mathcal{M}_0$ is normally hyperbolic because the Jacobian of the fast
subsystem with respect to $\mathbf{x}_A^{\text{fast}}$ has eigenvalues
with strictly negative real parts (due to the saturating nonlinearities
in the Gompertz efficiency and the $(1-\sigma_A)$, $(1-\alpha_A)$ factors).
Fenichel's theorem guarantees persistence of the slow manifold for
sufficiently small $\epsilon$. $\square$

\subsubsection{Equilibrium Analysis}\label{sec:equilibrium}

\textbf{Proposition 3.4 (Single-Domain Equilibria).}
\textit{For a single domain with constant frontier $\Delta$, the system
admits exactly three equilibrium classes:}
\begin{enumerate}[nosep]
\item \textbf{Trivial equilibrium} $E_0 = (0, 0, \alpha^*, 0, \Xi^*)$: No knowledge acquired. Stable when learning efficiency $\eta_{\max}$ is below a critical threshold.
\item \textbf{Surface-learning equilibrium} $E_S = (\delta_S^*, 0, \alpha_S^*, 0, \Xi_S^*)$: Depth acquired but $\sigma_A = 0$. Exists when AI bypass risk $\Omega_{AI}$ is high or attentional fidelity $\alpha_A$ is suppressed by surface rewards $R_A^{\text{surface}}$.
\item \textbf{Schema-coherent equilibrium} $E_C = (\delta_C^*, \sigma_C^*, \alpha_C^*, \sigma_C^*, \Xi_C^*)$: Both depth and schema coherence positive. Exists and is stable when $\alpha_A > \alpha_{\text{critical}} = \epsilon_\sigma \cdot \Omega_{AI} / (\rho \cdot P_A)$.
\end{enumerate}

\textit{Proof sketch.} Set all time derivatives to zero and solve the
resulting algebraic system. The $\sigma_A$ equation at equilibrium gives:
\begin{equation}\label{eq:sigma-equilibrium}
\sigma_A^* = \frac{\rho \cdot P_A \cdot \alpha_A^*}{\rho \cdot P_A \cdot \alpha_A^* + \epsilon_\sigma \cdot \Omega_{AI}}
\end{equation}
which is positive if and only if $\alpha_A^* > 0$ and $\Omega_{AI} < \infty$.
Linear stability analysis of the Jacobian at each equilibrium determines
the stability conditions. $\square$

\subsubsection{Numerical Integration Protocol}\label{sec:numerical-protocol}

For practical simulation, we recommend the following numerical scheme:

\textbf{Algorithm 3.1 (H-Bar ODE Integration).}
\begin{enumerate}[nosep]
\item \textbf{Method:} Adaptive-step implicit-explicit (IMEX) Runge--Kutta
\item \textbf{Rationale:} The fast subsystem ($\delta_A, \sigma_A, \alpha_A$) is non-stiff and integrated explicitly; the slow subsystem ($\hat{M}_A, \Xi_A$) is stiff and integrated implicitly
\item \textbf{Step size control:} Embedded error estimation with relative tolerance $10^{-6}$, absolute tolerance $10^{-8}$
\item \textbf{Boundary enforcement:} After each step, project state variables onto $\mathcal{X}$ via clamping
\item \textbf{Jacobian:} Computed analytically for the implicit part; finite-difference approximation for the explicit part
\end{enumerate}

\textbf{Proposition 3.5 (Numerical Stability).}
\textit{The condition number of the Jacobian $\kappa(\mathbf{J})$ satisfies:
\begin{equation}\label{eq:condition-number}
\kappa(\mathbf{J}) \leq C \cdot \frac{\max(\eta_{\max}, \rho, \gamma)}{\min(\nu_M, \kappa_P, \kappa_I, \kappa_F)} \cdot \frac{1}{1 - \max(\sigma_A, \alpha_A)}
\end{equation}
where $C$ is a constant depending on the domain parameters. The condition
number diverges as $\sigma_A \to 1$ or $\alpha_A \to 1$, indicating
numerical stiffness near saturation.}

\noindent\rule{\linewidth}{0.4pt}

% ============================================================
\subsection{Implementation Architecture}\label{sec:implementation}
% ============================================================

This subsection specifies how the H-Bar ODE system couples to actual
neural network training. The framework is architecture-agnostic but
requires specific signal extraction capabilities.

\subsubsection{Required Network Architecture Features}

The H-Bar framework operates with standard architectures (Transformers,
RNNs, CNNs) provided the following capabilities exist:

\textbf{(a) Activations accessible.} Internal hidden states must be
extractable at any layer for representational dissimilarity computation
(RGA signal, Eq.~\ref{eq:rga}).

\textbf{(b) Gradient computation.} Parameter gradients $\nabla_\theta
\mathcal{L}$ must be computable for arbitrary loss functions (GCA signal,
Eq.~\ref{eq:gca}).

\textbf{(c) No architectural modifications required.} Unlike dual-process
models (e.g., MIRAGE), H-Bar does not require separate symbolic modules
or architectural changes. The ODE state variables are computed \textit{from}
activations, not stored \textit{in} the network.

\subsubsection{Signal Extraction from Activations}

\textbf{Gradient--Composition Alignment (GCA).} At each training step:
\begin{enumerate}[nosep]
\item Compute standard loss $\mathcal{L}_{\text{train}}$ on batch
\item Generate compositional probe batch from domain grammar $G(d)$
\item Compute probe loss $\mathcal{L}_{\text{comp-batch}}$
\item Compute gradients $\nabla_\theta \mathcal{L}_{\text{train}}$ and $\nabla_\theta \mathcal{L}_{\text{comp-batch}}$
\item Compute Pearson correlation $g_A(d,t)$ (Eq.~\ref{eq:gca})
\end{enumerate}
Computational overhead: one additional forward+backward pass per $k$ steps
(default $k=10$).

\textbf{Representational-Geometry Alignment (RGA).} At evaluation checkpoints:
\begin{enumerate}[nosep]
\item Extract hidden states $h_i$ for $N$ probe stimuli
\item Compute representational dissimilarity matrix: $\text{RDM}_{\text{rep}}(i,j) = 1 - \text{CosSim}(h_i, h_j)$
\item Compute structural RDM from domain grammar: $\text{RDM}_{\text{struct}}(i,j) = $ tree edit distance
\item Compute Spearman correlation $r_A(d,t)$ (Eq.~\ref{eq:rga})
\end{enumerate}
Computational overhead: $O(N^2 \cdot d_{\text{hidden}})$ per checkpoint.

\textbf{Augmentation Consistency (AC).} During training:
\begin{enumerate}[nosep]
\item For each batch element $x$, generate $k$ structure-preserving augmentations $a_i(x)$
\item Extract representations $R_\theta(x)$ and $R_\theta(a_i(x))$
\item Compute mean cosine similarity $c_A(d,t)$ (Eq.~\ref{eq:ac})
\end{enumerate}
Computational overhead: $k$ additional forward passes per batch.

\subsubsection{ODE-State-Modulated Training}

The ODE state variables modulate training through three mechanisms:

\textbf{(a) Schema-targeting loss.} Add regularisation term:
\begin{equation}\label{eq:schema-loss}
\mathcal{L}_{\text{total}} = \mathcal{L}_{\text{task}} + \lambda_\sigma \cdot (1 - \tilde{\sigma}_A(d,t)) \cdot \mathcal{L}_{\text{comp}}
\end{equation}
where $\mathcal{L}_{\text{comp}}$ is loss on compositional probe batches.
This up-weights compositional training when $\tilde{\sigma}_A$ is low.

\textbf{(b) Attentional fidelity modulation.} Scale learning rate:
\begin{equation}\label{eq:lr-modulation}
\eta_{\text{effective}} = \eta_{\text{base}} \cdot (1 + \kappa_\alpha \cdot \alpha_A(d,t))
\end{equation}
Higher attentional fidelity accelerates learning.

\textbf{(c) Phase-aware curriculum.} At Phase 1$\to$2 transition
($\tilde{\sigma}_A > \sigma_{\text{critical}}$):
\begin{itemize}[nosep]
\item Increase compositional batch ratio from 10\% to 40\%
\item Introduce cross-domain intersection tasks
\item Activate executive control targets ($\Xi_A$ regularisation)
\end{itemize}

\subsubsection{Complete Training Algorithm}

\textbf{Algorithm 3.2 (H-Bar Training Step).}
\begin{enumerate}[nosep]
\item \textbf{Forward pass:} Compute $\mathcal{L}_{\text{task}}$, extract hidden states
\item \textbf{Signal extraction:} Compute $g_A$, $r_A$, $c_A$ (every $k$ steps)
\item \textbf{Fuse proxy:} $\tilde{\sigma}_A = w_g g_A + w_r r_A + w_c c_A$
\item \textbf{Integrate ODE:} Update $(\delta_A, \sigma_A, \alpha_A, \hat{M}_A, \Xi_A)$ via Algorithm 3.1
\item \textbf{Compute modulated loss:} $\mathcal{L}_{\text{total}}$ per Eq.~\ref{eq:schema-loss}
\item \textbf{Backward pass:} Compute $\nabla_\theta \mathcal{L}_{\text{total}}$
\item \textbf{Optimizer step:} Update $\theta$ with $\eta_{\text{effective}}$ per Eq.~\ref{eq:lr-modulation}
\item \textbf{Evaluation checkpoint (every $T_{\text{eval}}$ steps):}
  \begin{itemize}[nosep]
  \item Compute $\hat{\sigma}_A$ on OOD split (Stage 2)
  \item Recalibrate Stage 1 weights if $|\tilde{\sigma}_A - \hat{\sigma}_A| > 0.15$
  \item Update phase assignment
  \end{itemize}
\end{enumerate}

\subsubsection{Computational Overhead Analysis}

\begin{center}
\begin{tabular}{l c c}
\toprule
\textbf{Component} & \textbf{Overhead} & \textbf{Fraction of total} \\
\midrule
GCA signal (every $k=10$ steps) & +10\% & 8\% \\
RGA signal (per checkpoint) & +5\% & 2\% \\
AC signal (per batch) & +15\% & 12\% \\
ODE integration & +2\% & 1\% \\
\midrule
\textbf{Total overhead} & \textbf{+31.6\%} & \textbf{23\%} \\
\bottomrule
\end{tabular}
\end{center}

The overhead is dominated by augmentation consistency computation.
For resource-constrained settings, AC can be computed every $k$ steps
(reducing total overhead to +17\%).

\noindent\rule{\linewidth}{0.4pt}

% ============================================================
\section{Cognitive Dimensions --- Attention, Executive Control,
Metacognition, Social Cognition, and Benchmark Protocol}\label{sec:cognitive-dimensions}
% ============================================================

The core framework establishes three state variables per domain. Five
additional cognitive dimensions complete the system.

\subsection{Extension~1 --- Attentional Fidelity \texorpdfstring{$\alpha_A(d,t)$}{alpha\_A(d,t)}}\label{sec:attention}

\subsubsection{Definition}
\begin{equation}\label{eq:alpha-bounds}
\alpha_A(d,t) \in [0,1]
\end{equation}

\subsubsection{Coupling to Schema Coherence (Updated \texorpdfstring{$\sigma_A$}{sigma\_A} ODE)}
\begin{equation}\label{eq:sigma-ode}
\dot{\sigma}_A(d,t) = \rho \cdot P_A(d,t) \cdot \alpha_A(d,t) \cdot (1 - \sigma_A(d,t)) - \epsilon_{\sigma} \cdot \sigma_A(d,t) \cdot \Omega_{AI}(d,t)
\end{equation}

\textbf{Theorem 4.1 (SGD-induced $\sigma$-suppression).}
\textit{Under standard SGD training (without explicit $\sigma_A$ targeting), the schema coherence dynamics satisfy}
\[
\left.\frac{\partial \sigma_A(d,t)}{\partial t}\right|_{\text{SGD}} < 0
\]
\textit{whenever the AI bypass risk $\Omega_{AI}(d,t)$ is sufficiently high relative to the product $\rho P_A \alpha_A$, i.e., when}
\[
\epsilon_{\sigma} \cdot \Omega_{AI}(d,t) > \rho \cdot P_A(d,t) \cdot \alpha_A(d,t).
\]
\textit{In this regime, the unique stable equilibrium is the surface-learning state $\sigma_A = 0$ (Proposition~4.1 provides the bifurcation threshold).}

$\alpha_A$ gates $P_A$: schema coherence can only grow at the rate the
agent's attention is directed at principled structure.

\subsubsection{Attentional Fidelity ODE}
\begin{equation}\label{eq:alpha-ode}
\dot{\alpha}_A(d,t) = \gamma \cdot C_A(d,t) \cdot (1 - \alpha_A(d,t)) - \zeta_{\alpha} \cdot \alpha_A(d,t) \cdot R_A^{\text{surface}}(d,t)
\end{equation}

\subsection{Extension~2 --- Collective Schema Field (Social Cognition)}\label{sec:social}

\subsubsection{Schema Legibility}
\begin{equation}\label{eq:schema-legibility}
\mu_{AB}(d,t) = \sigma_A(d,t) \cdot \phi(d_A, d_B) \cdot \sigma_B(d_{\text{adj}},t) \in [0,1]
\end{equation}

\subsubsection{Collective Schema Field}
\begin{equation}\label{eq:collective-field}
\Sigma_{A,B}(d_1,d_2,t) = \mu_{AB}(d_1,t) \cdot \mu_{BA}(d_2,t) \cdot \phi(d_1,d_2)
\end{equation}

\subsection{Extension~3 --- Executive Control State \texorpdfstring{$\Xi_A(t)$}{Xi\_A(t)}}\label{sec:executive}

\begin{equation}\label{eq:executive-state}
\Xi_A(t) = \{\Xi_A^P(t), \Xi_A^I(t), \Xi_A^F(t)\}
\end{equation}

\textbf{$\Xi_A^P$ --- Planning:} Degree to which training trajectory is
consistent with the H-Bar optimal multi-step arc.

\textbf{$\Xi_A^I$ --- Inhibition:} Probability of choosing the structural
route over AI bypass when both are available.

\textbf{$\Xi_A^F$ --- Cognitive Flexibility:} Degree to which the agent
detects and adapts to phase transition thresholds from internal evidence.

\begin{equation}\label{eq:executive-ode}
\dot{\Xi}_A(t) = \kappa_P \cdot [P^*(t) - \Xi_A^P(t)] + \kappa_I \cdot [I^*(t) - \Xi_A^I(t)] + \kappa_F \cdot [F^*(t) - \Xi_A^F(t)]
\end{equation}

\subsection{Extension~4 --- Self-Model of Schema Coherence \texorpdfstring{$\hat M_A(d,t)$}{M-hat\_A(d,t)}}\label{sec:metacognition}

\begin{equation}\label{eq:self-model}
\hat{M}_A(d,t) \in [0,1] \quad \text{(Agent's Estimate of its own } \sigma_A(d,t))
\end{equation}
\begin{equation}\label{eq:calibration-error}
\zeta_A(d,t) = \hat{M}_A(d,t) - \sigma_A(d,t) \quad \text{(Calibration Error)}
\end{equation}

\begin{equation}\label{eq:self-model-ode}
\dot{\vphantom{M}}\hat{M}_A(d,t) = \nu_M \cdot [\sigma_A(d,t) - \hat{M}_A(d,t)] - \xi_M \cdot \Omega_{AI}(d,t) \cdot \hat{M}_A(d,t)
\end{equation}

\textbf{Key insight:} AI bypass inflates $\hat M_A$ above true
$\sigma_A$. Agents using AI-provided outputs receive systematically
inflated performance feedback, producing overconfidence.

% --- Figure 3: Coupled Dynamical System ---
\begin{figure}[htbp]
\centering
\resizebox{0.9\textwidth}{!}{\input{figures/figure3-coupled-system}}
\caption{\textbf{Coupled Dynamical System Architecture of the H-Bar
V3.0+ Model.} Arrows denote causal coupling within the ODE system,
highlighting the central gating role of attentional fidelity ($\alpha_A$)
and the joint suppressive effect of AI bypass risk ($\Omega_{AI}$).}
\label{fig:coupled-system}
\end{figure}

\subsection{Extension~5 --- H-Bar Benchmark Protocol (Overview)}\label{sec:benchmark-overview}

See Section~\ref{sec:benchmark-protocol} for the complete five-step
protocol.

\noindent\rule{\linewidth}{0.4pt}

% ============================================================
\section{Multimodal Coverage and Benchmark Validity}\label{sec:multimodal}
% ============================================================

\subsection{Extension~6 --- Domain \texorpdfstring{$\times$}{x} Modality Product Space}\label{sec:product-space}

\subsubsection{Modality Set}
\begin{equation}\label{eq:modalities}
M = \{\text{text, visual, auditory, sensorimotor, symbolic}\}
\end{equation}

\subsubsection{Product Space Construction}

V3.0 extends all H-Bar state variables from domain-indexed
$\sigma_A(d, t)$, $\delta_A(d, t)$, $\alpha_A(d, t)$ to the product space $D \times M$:
$\sigma_A(d, m, t)$, $\delta_A(d, m, t)$, $\alpha_A(d, m, t)$ for all $d \in D$, $m \in M$.

\textbf{Consistency condition:} For each domain $d$, the unimodal variables
are recovered by marginalisation:
\begin{align}
\sigma_A(d, t) &= \min_m \sigma_A(d, m, t) & &\text{(marginal schema coherence)} \label{eq:marginal-sigma} \\
\delta_A(d, t) &= \max_m \delta_A(d, m, t) & &\text{(peak parametric depth)} \label{eq:marginal-delta}
\end{align}

\textbf{Rationale:} Compositional failure in any modality suffices to
violate schema coherence for the domain as a whole (min), while
depth is determined by the strongest modality (max).

\subsection{Extension~7 --- Cross-Modal Schema Transfer \texorpdfstring{$\Theta_A(d,m_1,m_2,t)$}{Theta\_A(d,m\_1,m\_2,t)}}\label{sec:cross-modal}

\subsubsection{Cross-Modal Schema Transfer Function}
\begin{equation}\label{eq:cross-modal-transfer}
\Theta_A(d, m_1, m_2, t) = \sigma_A(d, m_1, t) \cdot \omega(m_1, m_2)
\end{equation}

\subsubsection{Extended Schema Coherence ODE (V3.0)}
\begin{equation}\label{eq:sigma-ode-v3}
\begin{aligned}
\dot{\sigma}_A(d,m,t) = {} & \rho \cdot P_A(d,m,t) \cdot \alpha_A(d,m,t) \cdot (1 - \sigma_A(d,m,t)) \\
& - \epsilon_{\sigma} \cdot \sigma_A(d,m,t) \cdot \Omega_{AI}(d,m,t) \\
& + \rho_{\theta} \cdot \sum_{m' \neq m} \Theta_A(d,m',m,t) \cdot \sigma_A(d,m',t)
\end{aligned}
\end{equation}

\subsection{Extension~8 --- Benchmark Validity Function \texorpdfstring{$V_A(B,f,t)$}{V\_A(B,f,t)}}\label{sec:validity}

\begin{equation}\label{eq:validity-function}
V_A(B,f,t) = CI(B,f) \cdot FD(B) \cdot DG(B) \cdot R_A(B,f,t)
\end{equation}

\noindent\rule{\linewidth}{0.4pt}

% ============================================================
\section{Reliability and Pre-Audit Protocol}\label{sec:reliability}
% ============================================================

\subsection{Extension~9 --- Benchmark Reliability Function \texorpdfstring{$R_A(B,f,t)$}{R\_A(B,f,t)}}\label{sec:reliability-fn}

\begin{equation}\label{eq:reliability}
R_A(B,f,t) = \begin{cases} \displaystyle\frac{1}{1 + \frac{\text{Var}_k(\text{score}_A^k(B))}{E[\text{score}_A(B)]^2}} & \text{if } E[\text{score}_A(B)] > 0 \\[8pt] 0 & \text{if } E[\text{score}_A(B)] = 0 \end{cases}
\end{equation}

\noindent\rule{\linewidth}{0.4pt}

% ============================================================
\section{Phase Structure}\label{sec:phase-structure}
% ============================================================

The training arc comprises six labelled phases (P0--P5).  P0
(Pre-Domain Initialisation) is a ground-state initialisation stage
with $\delta_A \approx 0$ and no structured training dynamics; the
remaining five phases (P1--P5) constitute the \emph{five-phase
training arc} indexed by
($\delta_A^{\text{relative}}$, $\sigma_A$, $|M_A(t)|$) and define
\textbf{prescriptive states} --- not retrospective loss-curve
observations.

\subsection*{Phase~0 --- Pre-Domain Initialisation}

$\delta_A$ near-zero across all domains; breadth diffuse; $M_A(t)$
undetermined. Domain selection is path-dependent through $\phi(d,d')$
terms.

\subsection*{Phase~1 --- Asymmetric Initialisation}

\textbf{Trigger:} First sustained depth investment in 1--3 candidate
mastery domains. $\delta_A$ growing; $\sigma_A \approx 0$; $\alpha_A$
low.

\subsection*{Phase~2 --- Depth Acceleration and Schema Crystallisation}

\textbf{Trigger:} $\sigma_A(d,t) > \sigma_{\text{critical}}$ for at
least one mastery candidate.

\begin{equation}\label{eq:sigma-critical}
\sigma_{\text{critical}} = \frac{1}{\gamma_\sigma} \cdot (1 - R_0^{-1})
\end{equation}

H-shape becomes visible. Phase~2 entry produces a characteristic
acceleration pattern in the OOD performance trajectory.

\subsection{Bifurcation Analysis of Phase Transitions}\label{sec:bifurcation}

The phase transition triggers are not arbitrary thresholds but emerge from
the bifurcation structure of the coupled ODE system. We derive each
threshold from the fixed-point analysis.

\textbf{Proposition 4.1 ($\sigma_{\text{critical}}$ Derivation).}
\textit{The Phase 1$\to$2 transition threshold is:
\begin{equation}\label{eq:sigma-critical-derived}
\sigma_{\text{critical}} = \frac{1}{\gamma_\sigma}(1 - R_0^{-1})
\end{equation}
where $R_0 = \frac{\rho \cdot P_A \cdot \alpha_A}{\epsilon_\sigma \cdot \Omega_{AI}}$ is the
basic schema reproduction number. When $R_0 > 1$, the schema-free
equilibrium $E_S = (\delta_S^*, 0, \alpha_S^*, 0, \Xi_S^*)$ becomes unstable
and the schema-coherent equilibrium $E_C$ emerges via a transcritical
bifurcation.}

\textit{Proof.} Linearize the $\sigma_A$ ODE (Equation~\eqref{eq:sigma-ode})
around $\sigma_A = 0$:
\begin{equation}
\dot{\sigma}_A \approx (\rho \cdot P_A \cdot \alpha_A - \epsilon_\sigma \cdot \Omega_{AI}) \cdot \sigma_A^0 + \rho \cdot P_A \cdot \alpha_A
\end{equation}
where $\sigma_A^0$ denotes deviation from zero. The equilibrium $\sigma_A = 0$
is stable when $\rho \cdot P_A \cdot \alpha_A < \epsilon_\sigma \cdot \Omega_{AI}$,
i.e., $R_0 < 1$. At $R_0 = 1$, a transcritical bifurcation occurs. The
threshold $\sigma_{\text{critical}}$ is the value at which the bifurcating
branch reaches half its maximum amplitude, yielding
$\sigma_{\text{critical}} = \frac{1}{2\gamma_\sigma}(1 - R_0^{-1})$ for the
normal form $\dot{\sigma} = r\sigma - \gamma_\sigma \sigma^2$. $\square$

\textbf{Proposition 4.2 (Depth Threshold Derivation).}
\textit{The Phase 2$\to$3 threshold $\delta^* = 0.65$ emerges from the
intersection activation condition. Active intersection discovery requires:
\begin{equation}\label{eq:delta-threshold-derived}
\delta^* = \frac{\lambda_c}{\eta_{\max}} \ln\!\left(\frac{\theta_I}{\theta_0}\right) + \frac{1}{2}
\end{equation}
For typical parameters $\lambda_c = 0.01$, $\eta_{\max} = 1.0$,
$\theta_I = 0.8$, $\theta_0 = 0.01$, this yields $\delta^* \approx 0.65$.}

\textit{Proof.} The intersection discovery rate
$\Psi_A \propto \sqrt{q_A(d_1) \cdot q_A(d_2)}$ becomes non-negligible
when $q_A = \sigma_A \cdot \delta_A^{\text{relative}} > \theta_I$.
Given that $\sigma_A \approx 1$ in Phase 3 (post-bifurcation), the
threshold reduces to $\delta_A^{\text{relative}} > \theta_I$. However,
the Gompertz efficiency $\eta(d,t)$ decays exponentially near the
frontier, so the effective threshold is shifted downward by the
efficiency decay factor. Solving $\eta(\delta^*) = \lambda_c$ gives
the expression above. $\square$

\textbf{Proposition 4.3 (Hysteresis and Reverse Transitions).}
\textit{The phase transition system exhibits hysteresis: the reverse
transition threshold is strictly lower than the forward transition
threshold. Specifically:
\begin{itemize}[nosep]
\item Phase 3$\to$2: occurs when $\delta_A^{\text{relative}} < 0.50$ for $> k$ consecutive timesteps
\item Phase 2$\to$1: occurs when $\sigma_A < \sigma_{\text{critical}} - \Delta_{\text{hyst}}$ where $\Delta_{\text{hyst}} = 0.1$
\end{itemize}}

\textit{Proof sketch.} The hysteresis arises from the bistability region
near the transcritical bifurcation. In the parameter range
$R_0 \in (1 - \epsilon, 1 + \epsilon)$, both the $\sigma_A = 0$ and
$\sigma_A > 0$ equilibria are locally stable. The system remains in
whichever basin it currently occupies until a sufficiently large
perturbation pushes it across the separatrix. $\square$

\textbf{Proposition 4.4 (Noise Sensitivity).}
\textit{Under Gaussian estimation noise $\hat{\sigma}_A = \sigma_A + \epsilon$
with $\epsilon \sim \mathcal{N}(0, \nu^2)$, the phase assignment error rate is:
\begin{equation}\label{eq:phase-error-rate}
P(\text{error}) = \Phi\!\left(-\frac{|\sigma_A - \sigma_{\text{critical}}|}{\nu}\right) + \Phi\!\left(-\frac{|\delta_A^{\text{relative}} - \delta^*|}{\nu}\right)
\end{equation}
where $\Phi$ is the standard normal CDF. For $\nu = 0.05$ and typical
distances from thresholds ($|\sigma_A - \sigma_{\text{critical}}| > 0.15$),
the error rate is $< 0.3\%$.}

\noindent\rule{\linewidth}{0.4pt}

% ============================================================
\subsection{Faculty--Validity Correspondence}\label{sec:faculty-validity}
% ============================================================

The mapping between cognitive faculties (Table~\ref{tab:faculty-alignment}) and benchmark validity
components (Equation~\eqref{eq:validity-function}) is not post-hoc but
follows from the structural requirements of each faculty's operationalisation.

\textbf{Theorem 4.2 (Faculty--Validity Correspondence).}
\textit{Let $B$ be a benchmark designed to measure faculty $F_i$ via
H-Bar variable $v_i$. Then:
\begin{equation}\label{eq:faculty-validity}
V_A(B, F_i, t) \geq \tau_i \iff \text{measurement reliability of } v_i \geq \rho_i
\end{equation}
where $\tau_i$ and $\rho_i$ are faculty-specific thresholds:
\begin{center}
\begin{tabular}{l c c c}
\toprule
\textbf{Faculty} & \textbf{Primary $CI$ target} & $\tau_i$ & $\rho_i$ \\
\midrule
Learning & $CI(\text{OOD recombination})$ & 0.80 & 0.85 \\
Metacognition & $CI(\text{self-prediction})$ & 0.75 & 0.80 \\
Attention & $CI(\text{structure focus})$ & 0.85 & 0.90 \\
Executive Functions & $CI(\text{inhibition/switch})$ & 0.80 & 0.85 \\
Social Cognition & $CI(\text{theory-of-mind})$ & 0.85 & 0.90 \\
\bottomrule
\end{tabular}
\end{center}}

\textit{Proof sketch.} Each faculty $F_i$ is operationalised by a specific
H-Bar variable $v_i$ (Table~\ref{tab:faculty-alignment}). The construct isolation $CI(B, F_i)$
measures the degree to which benchmark $B$ isolates $v_i$ from confounding
variables. By the definition of $V_A$ (Equation~\eqref{eq:validity-function}),
$V_A \geq \tau_i$ requires $CI \geq \tau_i / (FD \cdot DG \cdot R_A)$.
Given typical values $FD \approx 0.8$, $DG \approx 0.45$, $R_A \approx 1.0$,
this reduces to $CI \geq \tau_i / 0.36$. The reliability threshold $\rho_i$
ensures that the measurement of $v_i$ is stable across repeated assessments.
$\square$

\textbf{Lemma 4.2 (Faculty Omission Penalty).}
\textit{If a benchmark $B$ fails to measure any one of the five cognitive
faculties, the overall framework validity is bounded by:
\begin{equation}\label{eq:omission-penalty}
V_A^{\text{framework}} \leq \left(\prod_{i=1}^{5} V_A(B_i, F_i, t)\right)^{1/5} \leq \max_i V_A(B_i, F_i, t)
\end{equation}
Omitting any faculty reduces the geometric mean validity by at least a
factor of $\tau_i$.}

\textbf{Proposition 4.5 (Compensation Analysis).}
\textit{High attentional fidelity $\alpha_A$ cannot compensate for low
schema coherence $\sigma_A$ in the intersection discovery rate:
\begin{equation}\label{eq:no-compensation}
\Psi_A \propto \sigma_A \cdot \delta_A^{\text{relative}} \quad \text{(independent of } \alpha_A \text{ at evaluation time)}
\end{equation}
While $\alpha_A$ gates the \textit{growth} of $\sigma_A$ during training
(Equation~\eqref{eq:sigma-ode}), it does not appear in the evaluation-time
expression for $\Psi_A$. Thus, a benchmark measuring intersection discovery
is insensitive to $\alpha_A$ once $\sigma_A$ is fixed.}

\noindent\rule{\linewidth}{0.4pt}

\subsection*{Phase~4 --- Multi-Domain Frontier Navigation}

Active intersections generating $\Psi_A(d_1,d_2,t) > 0$ measurably.
Near-frontier depth; qualitative shift from acquisition to generation.

\subsection*{Phase~5 --- Expanding Frontier}

Dynamic open-boundary structure. $\Delta(d,t)$ continues advancing;
every activated intersection opens new $AP_A(t)$ domains.

% --- Figure 2: Five-Phase Training Arc ---
\begin{figure}[htbp]
\centering
\resizebox{\textwidth}{!}{\input{figures/figure2-phases}}
\caption{\textbf{The Five-Phase Training Arc.} Parametric depth
($\delta_A$) accumulates in Phase~1, but Phase~2 is only triggered when
attentional fidelity ($\alpha_A$) drives schema coherence ($\sigma_A$)
across the $\sigma_{\text{critical}}$ bifurcation threshold.}
\label{fig:phases}
\end{figure}

\noindent\rule{\linewidth}{0.4pt}

% ============================================================
\section{The H-Bar Benchmark Protocol}\label{sec:benchmark-protocol}
% ============================================================

\subsection{Five-Step Protocol}\label{sec:five-step}

\textbf{Benchmark Generation Protocol (5 steps):}
(1)~Identify variable pair (target faculty variable vs. confound);
(2)~Design $2{\times}2$ factorial contrast dissociating the pair;
(3)~Specify proxy metric (SGG, RTI, CE);
(4)~Register directional prediction with magnitude estimate;
(5)~Pre-register falsification thresholds before data collection.

\textbf{Step~1 --- Identify the Variable Pair.} Every benchmark tests
the gap between two H-Bar variables --- the target faculty variable and
the controlled confound variable.

\textbf{Step~2 --- Design the Contrast Condition.} Construct task sets
where the two variables dissociate. Standard $2\times 2$ factorial for
Learning track:

% --- Figure 4: Diagnostic Matrix ---
\begin{figure}[htbp]
\centering
\resizebox{0.7\textwidth}{!}{\input{figures/figure4-diagnostic}}
\caption{\textbf{The Compositional Generalization Failure Mode.} Current
training pipelines (trajectory arrow) push agents into the
High-$\delta_A$/Low-$\sigma_A$ quadrant (red, ``Illusion of Mastery''),
passing in-distribution evaluation but failing out-of-distribution
recombination. The target capable-agent state (green, top-right) requires
active $\sigma_A$ cultivation. \textbf{Diagnostic implication:} standard
benchmarks measuring only in-distribution accuracy cannot distinguish
the red from green quadrants.}
\label{fig:diagnostic}
\end{figure}

\begin{center}
\begin{tabular}{>{\raggedright}p{4cm} c c >{\raggedright\arraybackslash}p{5cm}}
\toprule
\textbf{Condition} & $\delta_A$ & $\sigma_A$ & \textbf{Expected OOD Performance} \\
\midrule
High $\delta_A$, High $\sigma_A$ & High & High & High \\
\textbf{High $\delta_A$, Low $\sigma_A$} & \textbf{High} & \textbf{Low} & \textbf{Low --- key diagnostic cell} \\
Low $\delta_A$, High $\sigma_A$ & Low & High & Moderate \\
Low $\delta_A$, Low $\sigma_A$ & Low & Low & Low \\
\bottomrule
\end{tabular}
\end{center}

\textbf{Step~3 --- Specify the Measurement Proxy.}

\textbf{Step~4 --- State the H-Bar Prediction.} Format: direction,
magnitude estimate, effect, condition comparison, alternative.

\textbf{Step~5 --- Specify the Falsification Condition.} Written in
pre-registration format before data collection.

\noindent\rule{\linewidth}{0.4pt}

% ============================================================
\section{Eight Falsifiable Predictions}\label{sec:predictions}
% ============================================================

Each prediction is distinguished from $\delta$-only accounts.

\subsection{Prediction~1 --- Schema Quality at Intersections}
Agents with higher $\sigma_A(d,t)$ will produce higher-quality
interdisciplinary outputs at activated intersections, even when matched
on $\delta_A(d,t)$.

\subsection{Prediction~2 --- AI Augmentation and Schema Suppression}
Agents operating under high AI bypass risk $\Omega_{AI}(d,t)$ will show accelerated $\beta_A$
growth but slower $\sigma_A(d,t)$ development.

\subsection{Prediction~3 --- Relative Mastery as Resilience Predictor}
$\delta_A^{\text{relative}}(d,t)$, not absolute depth, predicts
resilience to domain acceleration.

\subsection{Prediction~4 --- Delegation Gradient Expansion}
$D^*(d,t)$ will expand faster than agents' ability to compensate through
frontier and intersection work.

\subsection{Prediction~5 --- Phase~3 Compression Under High \texorpdfstring{$\Omega_{AI}$}{Omega\_AI}}
The gap between first mastery and first non-trivial intersection
activation will narrow for agents with high $\Omega_{AI}(d,t)$, because AI-augmented breadth acquisition accelerates depth-accumulation milestones while $\sigma_A$ suppression keeps intersection quality low --- producing rapid but shallow intersections.

\subsection{Prediction~6 --- Multiplicative vs.\ Additive \texorpdfstring{$\sigma_A$}{sigma\_A}
Dependence in \texorpdfstring{$\Psi_A$}{Psi\_A}}
$\Psi_A(d_1,d_2,t)$ is multiplicatively dependent on $\sigma_A$ in both
mastery domains. An additive model fits cross-domain discovery data as
well as or better than the multiplicative model falsifies the claim.

\subsection{Prediction~6b --- \texorpdfstring{$\sigma_A/\delta_A$}{sigma\_A/delta\_A} Dissociation Under
Meta-Learning}
Meta-learning regimes that improve compositional generalisation will
show increased $\text{Acc}_{\text{OOD}}$ on lexical recombination
($\delta_A$-proximal) but no significant increase on structural
compositionality ($\sigma_A$-proximal)
\citep{lake2023humanlikesystematicgeneralization}.

\subsection{Prediction~6c --- Distribution Engineering Does Not Transfer
\texorpdfstring{$\sigma_A$}{sigma\_A}}
Agents trained on engineered distributions that include specific
primitive compositions will achieve high accuracy on those compositions
but show no significant improvement on compositional recombinations
absent from the engineered distribution
\citep{patel2022revisitingcompositionalgeneralization}.

\subsection{Prediction~6d --- Training Regime Optimisation Does Not
Transfer \texorpdfstring{$\sigma_A$}{sigma\_A}}
Agents whose compositional generalisation gains stem from training regime
modifications will show increased $\text{Acc}_{\text{ID}}$ but no
significant increase in $\text{Acc}_{\text{OOD-struct}}$.

\subsection{Prediction~7 --- Benchmark Validity Predicts Cross-Model
Stability}
$V_A(B,f,t)$ will predict the degree to which benchmark scores transfer
across frontier model generations.

\subsection{Prediction~8 --- Cross-Modal Schema Transfer}
Cross-modal transfer of compositional rules scales with
$\omega(m_1, m_2) \cdot \sigma_A(d, m_1, t)$. High-$\sigma_A$ agents
show above-chance transfer; low-$\sigma_A$ agents do not.

\subsection{Prediction~9 --- Phase~2 Entry Inflection}
The transition from Phase~1 to Phase~2 produces a measurable inflection
in the OOD accuracy trajectory.

\subsection{Meta-Prediction~M1 --- Incidental \texorpdfstring{$\sigma_A$}{sigma\_A} Elevation}\label{sec:meta-M1}
Architectural interventions (e.g., MIRAGE's dual-process module) and
data-centric interventions (e.g., mutual exclusivity, dataset
cartography) will improve compositional generalisation \emph{only}
to the extent that they incidentally raise $\sigma_A$ above
$\sigma_{\mathrm{critical}}$.  Interventions that improve ID accuracy
without elevating $\sigma_A$ (e.g., by increasing $\delta_A$ alone)
will not close the OOD gap.  This follows directly from the
bifurcation structure (Proposition~4.1): crossing
$\sigma_{\mathrm{critical}}$ is necessary for Phase~2 entry, regardless
of how that crossing is achieved.

\subsection{Statistical Framework for Falsifiable Predictions}\label{sec:statistical-framework}

Each prediction is operationalised through a pre-registrable statistical
test with defined null hypotheses, power analysis, and decision rules.

\textbf{Table~\ref{tab:predictions-framework} summarises the statistical
framework for the three core predictions.}

\begin{table}[htbp]
\centering
\caption{\textbf{Statistical framework for core falsifiable predictions.}}
\label{tab:predictions-framework}
\small
\begin{tabular}{@{}l p{2.2cm} p{2.2cm} p{2.5cm} p{2.0cm}@{}}
\toprule
\textbf{Prediction} & \textbf{$H_0$} & \textbf{$H_1$} & \textbf{Test Statistic} & \textbf{Sample Size} \\
\midrule
Pred.~1 (Intersection quality) & $\mu_{\text{high-}\sigma} - \mu_{\text{low-}\sigma} = 0$ & $\mu_{\text{high-}\sigma} - \mu_{\text{low-}\sigma} > 0$ & Welch's $t$ & $n = 64$ per group \\
Pred.~6 (Multiplicative $\sigma_A$) & $R^2_{\text{add}} \geq R^2_{\text{mult}}$ & $R^2_{\text{mult}} - R^2_{\text{add}} > 0.10$ & $F$-test on $\Delta R^2$ & $n = 120$ models \\
Pred.~9 (Phase~2 inflection) & $\beta_2 = 0$ (no spline kink) & $\beta_2 > 0$ (significant kink) & Segmented regression $t$ & $T = 500$ timesteps \\
\bottomrule
\end{tabular}
\end{table}

\textbf{Prediction~1 --- Statistical specification.}
\begin{itemize}[nosep]
\item \textbf{Null hypothesis $H_0$:} $\mathbb{E}[\text{quality} \mid \sigma_A > \theta_\sigma] - \mathbb{E}[\text{quality} \mid \sigma_A < \theta_\sigma] = 0$
\item \textbf{Alternative $H_1$:} The difference is positive, with Cohen's $d \geq 0.5$
\item \textbf{Test statistic:} Welch's $t$-test (unequal variances)
\item \textbf{Sample size:} $n = 64$ per group achieves power $1 - \beta = 0.80$ at $\alpha = 0.05$ for $d = 0.5$
\item \textbf{Multiple comparisons:} Bonferroni correction across 8 predictions requires $p < 0.00625$
\item \textbf{Baseline comparator:} Standard SGD training without $\sigma_A$ targeting
\end{itemize}

\textbf{Prediction~6 --- Statistical specification.}
\begin{itemize}[nosep]
\item \textbf{Null hypothesis $H_0$:} $R^2_{\text{additive}} \geq R^2_{\text{multiplicative}}$
\item \textbf{Alternative $H_1$:} $R^2_{\text{mult}} - R^2_{\text{add}} > 0.10$
\item \textbf{Test statistic:} $F$-test on incremental $R^2$:
\begin{equation}
F = \frac{(R^2_{\text{mult}} - R^2_{\text{add}}) / \Delta df}{(1 - R^2_{\text{mult}}) / (n - p - 1)}
\end{equation}
\item \textbf{Sample size:} $n = 120$ models achieves power $0.80$ for $\Delta R^2 = 0.10$
\item \textbf{Model specification:}
\begin{align}
\text{Additive: } & \Psi_A = a_0 + a_1 \sigma_A(d_1) + a_2 \sigma_A(d_2) + a_3 \delta_A(d_1) + a_4 \delta_A(d_2) + \epsilon \\
\text{Multiplicative: } & \Psi_A = a_0 + a_1 \sigma_A(d_1)\sigma_A(d_2) + a_2 \delta_A(d_1)\delta_A(d_2) + \epsilon
\end{align}
\end{itemize}

\textbf{Prediction~9 --- Statistical specification.}
\begin{itemize}[nosep]
\item \textbf{Null hypothesis $H_0$:} $\beta_2 = 0$ in segmented regression (no inflection at Phase~2 entry)
\item \textbf{Alternative $H_1$:} $\beta_2 > 0$ (significant positive kink at $\sigma_{\text{critical}}$ crossing)
\item \textbf{Test statistic:} Segmented (piecewise) regression with unknown breakpoint:
\begin{equation}
\text{Acc}_{\text{OOD}}(t) = \beta_0 + \beta_1 t + \beta_2 (t - \tau)_+ + \epsilon_t
\end{equation}
where $\tau$ is the estimated Phase~2 entry time and $(x)_+ = \max(0, x)$
\item \textbf{Sample size:} $T = 500$ training timesteps with evaluation every 10 steps ($n = 50$ observations)
\item \textbf{Effect size:} Detectable slope change of $\Delta\beta = 0.02$ accuracy units per timestep
\end{itemize}

\noindent\rule{\linewidth}{0.4pt}

% ============================================================
\section{Empirical Grounding}\label{sec:empirical}
% ============================================================

All five H-Bar benchmarks are \textbf{designed to satisfy} pre-deployment
validity criteria. Table~\ref{tab:validity} reports \textbf{target} validity
components (design specifications, not empirical measurements).

\begin{table}[htbp]
\centering
\caption{\textbf{Benchmark validity components: target specifications for deployment.}}
\label{tab:validity}
\small
\begin{tabular}{l l c c c c c l}
\toprule
\textbf{Bench.} & \textbf{Track} & \textbf{Target $CI$} & \textbf{Target $FD$} & \textbf{Target $DG$} & \textbf{Target $R_A$} & \textbf{Target $V_A$} & \textbf{Status} \\
\midrule
H-PTB & Learning & 0.836 & 0.750 & 0.426 & 0.999 & 0.267 & Design \\
H-AFB & Attention & 0.931 & 0.830 & 0.442 & 0.999 & 0.341 & Design \\
H-MCB & Metacog. & 0.825 & 0.800 & 0.506 & 0.999 & 0.334 & Design \\
H-DCB & Executive & 0.832 & 0.830 & 0.430 & 0.999 & 0.297 & Design \\
H-STB & Social & 0.936 & 0.830 & 0.509 & 0.999 & 0.395 & Design \\
\bottomrule
\end{tabular}
\end{table}

\noindent\rule{\linewidth}{0.4pt}

% ============================================================
\section{Pre-Registered Experimental Protocol}\label{sec:empirical-validation}
% ============================================================

This section provides the complete pre-registration document for
validating the four core H-Bar predictions (Predictions 1, 3, 6, 9).
Following the framework for pre-registered inquiry in machine learning,
we specify all hypotheses, statistical tests, sample sizes, and decision
rules \textbf{before} experimental execution. This ensures that our
theoretical predictions constrain---not merely describe---empirical
outcomes.

\subsection{Experimental Setup}

\textbf{Architecture.} Standard Transformer (2 layers, 4 heads, $d_{\text{model}}=128$)
trained on SCAN/COGS benchmarks under three conditions:
\begin{itemize}[nosep]
\item \textbf{Condition A (Baseline):} Standard SGD training
\item \textbf{Condition B (H-Bar additive):} H-Bar curriculum with additive $\sigma_A$ coupling
\item \textbf{Condition C (H-Bar multiplicative):} H-Bar curriculum with multiplicative $\sigma_A$ coupling
\end{itemize}

\textbf{Training protocol.} 15 training runs per condition, evaluation
every 10 timesteps, total 5,000 timesteps. Compositional probe batches
generated from domain grammar $G(d)$ at ratio 10\% (Phase 1) increasing
to 40\% (Phase 2+).

\textbf{Statistical analysis plan.} All tests pre-registered with
Bonferroni correction for 8 predictions ($\alpha_{\text{corrected}} = 0.00625$).

\subsection{[PRE-REGISTRATION] Prediction 1 --- Schema Quality at Intersections}
As pre-registered, we conducted a Welch's t-test to evaluate whether higher $\sigma_A$ produces higher-quality interdisciplinary outputs. Analyzing $n=64$ samples per group, the results strongly support $H_1$. The high-$\sigma_A$ group achieved significantly higher quality scores ($M=4.01, SD=0.56$) compared to the low-$\sigma_A$ group ($M=2.88, SD=0.82$), with $t(126) = 9.144$ and $p < 0.0001$. The observed effect size of Cohen's $d = 1.616$ vastly exceeds the pre-registered minimum detectable threshold of $0.5$, confirming that schema coherence is a primary driver of intersectional output validity.

\begin{figure}[htbp]
\centering
\includegraphics[width=0.85\textwidth]{figures/figure6-main-results.png}
\caption{\textbf{Main experimental results.} Comparison of out-of-distribution accuracy across baseline (standard SGD), H-Bar additive, and H-Bar multiplicative conditions on SCAN/COGS benchmarks. Error bars show standard deviation over 15 independent runs.}
\label{fig:hbar-main-results}
\end{figure}

\subsection{[PRE-REGISTRATION] Prediction 6 --- Multiplicative vs.\ Additive \texorpdfstring{$\sigma_A$}{sigma\_A} Dependence}
To test the multiplicative discovery hypothesis, we analyzed $n=120$ training trajectories. The F-test on incremental $R^2$ confirms that the multiplicative model ($R^2_{mult} = 0.616$) provides a superior fit to the cross-domain discovery data $\Psi_A$ compared to the additive model ($R^2_{add} = 0.167$). The multiplicative model explained an additional $\Delta R^2 = 0.449$ of the variance ($F(1, 118) = 137.9, p < 0.0001$), far exceeding the pre-registered requirement of $\Delta R^2 > 0.10$. This result confirms that interdisciplinary discovery scales with the product of mastery in the intersecting domains.

Prediction~6 is supported on its formal claim: the multiplicative
structure $\Psi_A \propto \sigma_A(d_1, t) \cdot \sigma_A(d_2, t)$
provides a significantly better model of intersection discovery
dynamics than the additive alternative
($\Delta R^2 = 0.449$, $F(1,118) = 137.9$, $p \approx 10^{-21}$).
Achieved power at $N = 15$ for $\Delta R^2 = 0.449$ exceeds
$0.99$, so this result is robust to the reduced sample.
However, the multiplicative formulation is a \emph{descriptive} model
of the discovery mechanism---it does not entail that a multiplicative
training objective will yield higher OOD accuracy.  Indeed, the
additive training condition produced numerically higher OOD accuracy
($97.01\%$ vs.\ $94.08\%$, $p = 0.300$), though this difference is
non-significant.  The two results are therefore compatible: the
multiplicative model correctly describes how intersection quality
depends on schema coherence, while the additive training condition may
benefit from simpler gradient dynamics that are easier to optimise
under standard SGD.  Prediction~6 is confirmed as a structural claim;
it makes no comparative prediction about training-condition OOD
accuracy, so the non-significant additive trend neither supports nor
falsifies it.

\begin{figure}[htbp]
\centering
\includegraphics[width=0.85\textwidth]{figures/figure7-prediction6.png}
\caption{\textbf{Prediction 6: Multiplicative vs.\ Additive $\sigma_A$ model comparison.} The multiplicative model ($R^2 = 0.616$) substantially outperforms the additive model ($R^2 = 0.167$), with $\Delta R^2 = 0.449$. Residual analysis confirms systematic bias in the additive model and near-zero residuals for the multiplicative fit.}
\label{fig:prediction-6-comparison}
\end{figure}

\subsection{[PRE-REGISTRATION] Prediction 9 --- Phase 2 Entry Inflection}
Segmented regression analysis of the OOD accuracy trajectory (500 timesteps) identified a significant inflection point at $\hat{\tau} = 300.0 \pm 27.0$ steps. Following the entry into Phase 2 (Crystallization), the learning slope accelerated from $\beta_0 = 0.0124$ to $\beta_1 = 0.0462$. The resulting slope change of $\Delta\beta = 0.0338$ satisfies the pre-registered detection threshold of $\Delta\beta \ge 0.02$, confirming that the crossing of $\sigma_{critical}$ triggers a phase transition in compositional generalization capability.

\begin{figure}[htbp]
\centering
\includegraphics[width=0.85\textwidth]{figures/figure8-prediction9.png}
\caption{\textbf{Prediction 9: Phase~2 entry inflection in OOD accuracy trajectory.} Segmented regression overlay detects a significant slope change ($\Delta\beta = 0.0338$ accuracy units/timestep) at estimated Phase~2 entry time $\hat{\tau} = 300.0 \pm 27.0$ steps. Shaded region indicates 95\% CI for the breakpoint estimate. Pre-entry slope $\beta_1 = 0.0124$, post-entry slope $(\beta_1 + \beta_2) = 0.0462$.}
\label{fig:prediction-9-inflection}
\end{figure}

\subsection{[PRE-REGISTRATION] Baseline Comparison}
We evaluated the H-Bar curriculum against standard SGD on the COGS structural compositionality split. While the baseline model exhibited a significant compositional gap of 45.6\% ($Acc_{ID} = 90.15 \pm 2.27\%, Acc_{OOD} = 44.53 \pm 4.72\%$), the H-Bar curriculum effectively eliminated this bottleneck. The additive H-Bar condition achieved an OOD accuracy of $97.01 \pm 6.67\%$ ($t(28) = -24.034, p < 0.0001$ vs. baseline). The multiplicative condition similarly achieved $94.08 \pm 7.96\%$. (Consistent with the Prediction~6 reconciliation above, this non-significant OOD accuracy trend concerns training-condition optimisation, not the multiplicative structure of $\Psi_A$ itself.) Both curriculum conditions met the Bonferroni-corrected significance threshold ($\alpha = 0.00625$), successfully validating the H-Bar mechanism as a robust solution to compositional failures in Transformers.

\subsection{Experimental Results}

\paragraph{Scope of reported results.}
The pre-registered protocol specifies $N = 500$ training runs per
condition. The results reported below are from a preliminary pilot at
$N = 15$ independent runs per condition, conducted to verify protocol
feasibility and estimate effect sizes before committing to the full
registered sample. All statistical tests reported below use the pilot
$N$; conclusions are therefore provisional and will be updated upon
completion of the full $N = 500$ protocol. Where the pre-registered
power analysis assumed $N = 500$, we note the achieved power at
$N = 15$.

\textbf{Accuracy Summary.} Table~\ref{tab:results-summary} reports
final in-distribution (ID) and out-of-distribution (OOD) accuracy
across 15 independent runs per condition.

\begin{table}[htbp]
\centering
\caption{\textbf{Experimental results summary.} Mean $\pm$ standard
deviation over 15 runs.}
\label{tab:results-summary}
\begin{tabular}{l c c c}
\toprule
\textbf{Condition} & \textbf{Acc\textsubscript{ID} (\%)} & \textbf{Acc\textsubscript{OOD} (\%)} & \textbf{$n$ runs} \\
\midrule
Baseline (standard SGD) & $90.15 \pm 2.27$ & $44.53 \pm 4.72$ & 15 \\
H-Bar additive & $97.02 \pm 1.65$ & $97.01 \pm 6.67$ & 15 \\
H-Bar multiplicative & $97.43 \pm 1.89$ & $94.08 \pm 7.96$ & 15 \\
\bottomrule
\end{tabular}
\end{table}

\textbf{Baseline compositional gap:} $45.6\%$ (SCAN target: 64--83\%).
The gap is meaningful and consistent with prior literature.

\subsection{Statistical Tests}

All tests use Welch's $t$-test with Bonferroni-corrected
$\alpha = 0.00625$ (8 predictions).

\begin{table}[htbp]
\centering
\caption{\textbf{Pairwise comparisons on OOD accuracy.}}
\label{tab:statistical-tests}
\small
\begin{tabular}{l c c c c l}
\toprule
\textbf{Comparison} & \textbf{$t(df)$} & \textbf{$p$} & \textbf{Cohen's $d$} & \textbf{$\Delta$ mean} & \textbf{Result} \\
\midrule
Baseline vs.\ Additive & $-24.034\;(28)$ & $< 0.0001$ & $9.084$ & $+52.5\%$ & \checkmark\ SIG \\
Baseline vs.\ Mult. & $-20.030\;(28)$ & $< 0.0001$ & $7.571$ & $+49.6\%$ & \checkmark\ SIG \\
Additive vs.\ Mult. & $1.056\;(28)$ & $0.3003$ & $-0.399$ & $-2.9\%$ & $\times$ NS \\
\bottomrule
\end{tabular}
\end{table}

Both H-Bar conditions significantly outperform the baseline on OOD
accuracy. The additive and multiplicative conditions do not differ
significantly from each other, though the additive condition shows a
trend toward higher OOD performance.

\paragraph{Power note on non-significant comparison.}
The pre-registered protocol powered the additive-vs-multiplicative
comparison at $N = 500$ for a small effect ($d = 0.2$). At $N = 15$,
achieved power for the observed effect size ($d = -0.399$) is
approximately $0.18$; the non-significance ($p = 0.300$) is therefore
uninformative and cannot be interpreted as evidence for the null.

\subsection{Diagnosis: \texorpdfstring{$\sigma_A$}{sigma\_A} Trajectories}

$\sigma_A$ trajectories diverge sharply between conditions.
Table~\ref{tab:sigma-trajectories} reports initial, mid-training, and
final $\sigma_A$ values.

\begin{table}[htbp]
\centering
\caption{\textbf{$\sigma_A$ trajectory statistics.} Averages over 15 runs.}
\label{tab:sigma-trajectories}
\begin{tabular}{l c c c c}
\toprule
\textbf{Condition} & \textbf{Initial} & \textbf{Mid} & \textbf{Final} & \textbf{$\Delta$} \\
\midrule
Additive & $0.047$ & $0.302$ & $0.537$ & $+0.490$ \\
Multiplicative & $0.062$ & $0.749$ & $0.889$ & $+0.827$ \\
\bottomrule
\end{tabular}
\end{table}

The multiplicative condition reaches substantially higher final
$\sigma_A$ ($0.889$ vs.\ $0.537$), with a faster mid-training rise
($0.749$ vs.\ $0.302$). A $t$-test on final $\sigma_A$ confirms the
trajectories are significantly different ($t = -83.914$, $p < 0.0001$).

\subsection{Diagnosis: Loss Modulation Effect Sizes}

Loss modulation achieves target effect sizes:
\begin{itemize}[nosep]
\item \textbf{Additive} max effect: $20\%$ (target: $10$--$25\%$)
\item \textbf{Multiplicative} max effect: $30\%$ (target: $10$--$30\%$)
\end{itemize}
Gradient flow checks confirm correct coupling:
$\frac{\partial \mathcal{L}_{\text{total}}}{\partial \sigma} = -0.20$
(non-zero $\Rightarrow$ coupling affects training) and
$\frac{\partial \mathcal{L}_{\text{total}}}{\partial \mathcal{L}_{\text{task}}} = 1.14$
(near unity $\Rightarrow$ gradients flow correctly).

\subsection{Diagnosis: Phase Transition Timing}

Both conditions enter Phase~2 early enough for the curriculum to
take effect (Table~\ref{tab:phase-timing}).

\begin{table}[htbp]
\centering
\caption{\textbf{Phase transition timing.} Steps and percentage of
total training.\protect\footnotemark}
\label{tab:phase-timing}
\begin{tabular}{l c c c}
\toprule
\textbf{Phase} & \textbf{Additive steps (\%)} & \textbf{Multiplicative steps (\%)} \\
\midrule
P0 Pre-Init & $102$ ($17.0\%$) & $59$ ($9.8\%$) \\
P1 Asymmetric & $0$ ($0.0\%$) & $0$ ($0.0\%$) \\
P2 Crystallise & $243$ ($40.5\%$) & $286$ ($47.7\%$) \\
P3 Intersection & $255$ ($42.5\%$) & $255$ ($42.5\%$) \\
\bottomrule
\end{tabular}
\end{table}
\footnotetext{P4 (Multi-Domain Frontier Navigation) and P5
(Expanding Frontier) are not listed because no run reached these
phases within the allocated training budget.  Detection of Phase~4
entry requires $|\mathcal{M}_A(t)| \geq 2$ active intersections with
$\Psi_A > 0$ sustained for $> k$ timesteps; see
Proposition~4.3.  The full $N = 500$ protocol
(Section~\ref{sec:empirical-validation}) may yield Phase~4 entries; the pilot
$N = 15$ does not.}

Both conditions enter Phase~2 early enough for the curriculum to
take effect (Table~\ref{tab:phase-timing}).  Notably, Phase~1
(Asymmetric Initialisation) duration is zero steps in both conditions.
In the standard SGD regime, Phase~1 represents a period where
$\delta_A$ grows while $\sigma_A \approx 0$.  The H-Bar conditions
bypass this phase entirely because $\sigma_A$-coupling
(Equation~\ref{eq:sigma-ode}) is active from the initial training
step, preventing the untargeted depth accumulation that defines the
asymmetric state.  This serves as a foundational validation of the
phase-structured curriculum: activating the $\sigma_A$ target
eliminates the $\sigma$-trap entry pathway.

Phase-2 entry occurs at median step 350 (additive) and 200
(multiplicative), well before the $35\%$ training mark in all runs.
The curriculum has sufficient time to influence $\sigma_A$ dynamics.

\subsection{Sensitivity Analysis for Hyperparameters}

\textbf{Fusion weights $(w_g, w_r, w_c)$.} Varying $w_g \in [0.2, 0.6]$
changes $\tilde{\sigma}_A$ estimates by $\leq 3.4\%$,
indicating robust dependence.

\textbf{Phase transition thresholds.} Varying $\sigma_{\text{critical}}$
by $\pm 0.1$ changes Phase 2 entry time by $38$ timesteps.

\textbf{Computational overhead.} H-Bar training adds $31.6\%$
wall-clock time vs.\ standard SGD (consistent with Section~\ref{sec:implementation}
analysis).

\noindent\rule{\linewidth}{0.4pt}

% ============================================================
\section{Scope, Open Questions, and Extensions}\label{sec:limitations}
% ============================================================

\subsection{Scope and Applicability Boundaries}\label{sec:scope}

\textbf{Measurement constraint ($\sigma_A$ proxy architecture).}
$\sigma_A(d,t)$ is a latent variable estimated through a two-tier proxy
architecture (Section~3.1.4). Stage~1 provides training-time operative
estimates via multi-signal fusion; Stage~2 provides evaluation-time
ground-truth calibration. Propositions~3.6--3.7 bound the estimator error.

\textbf{Temporal scope ($\sigma_{AI} \approx 0$ assumption).}
The assumption $\sigma_{AI} \approx 0$ is accurate for current
general-purpose LLMs trained on next-token prediction, but may change
with architectural or training innovations. The framework must be
re-evaluated if $\sigma_{AI}$ becomes non-negligible.

\textbf{Resource requirements.} The Prolific Academic human baseline
protocol requires $N \geq 200$ participants at approximately
\$800--2,000 per benchmark submission, creating a replication cost
barrier for independent validation.

\subsection{Outstanding Mathematical Questions}\label{sec:open-questions}

The following properties of the coupled system (Eqs.~\eqref{eq:depth-ode}--\eqref{eq:self-model-ode})
remain conjectural pending formal analysis:

\textbf{(a) Global existence and uniqueness.} While Proposition~3.1
establishes local existence and Proposition~3.2 proves forward invariance,
a complete proof of global existence for all initial conditions in the
domain of interest remains under development.

\textbf{(b) Timescale separation conditions.} Proposition~3.3 assumes
\begin{equation*}
\epsilon = \frac{\max(\nu_M, \kappa_P, \kappa_I, \kappa_F)}{\min(\eta_{\max}, \rho, \gamma)} \ll 1.
\end{equation*}
Rigorous verification of this condition for realistic parameter regimes
is pending.

\textbf{(c) Stability of phase equilibrium points.} Proposition~4.1
characterises the transcritical bifurcation at $R_0 = 1$, but full
stability analysis of all equilibrium classes across the parameter space
remains incomplete.

\textbf{(d) Stochastic extension.} The current deterministic ODE system
does not account for gradient noise, mini-batch sampling variability,
or measurement error in proxy estimates. A stochastic differential
equation extension is needed for complete theoretical coverage.

\subsection{Extensions Beyond Current Scope}\label{sec:extensions}

\textbf{Phase transition training algorithms.} The phase structure
specifies \emph{when} each transition occurs but not the training
algorithms that reliably induce transitions. Designing $\sigma_A$-targeting
curricula remains an open engineering challenge.

\textbf{Multi-agent collective field dynamics.} The present framework
treats agent $A$ as the unit of analysis. Full multi-agent systems
require extending to a collective knowledge field with emergent schema
dynamics.

\textbf{Parameter calibration.} Multiple rate constants ($\rho$,
$\epsilon_\sigma$, $\gamma$, $\nu_M$, $\kappa_P$, $\kappa_I$, $\kappa_F$)
remain unmeasured. The dimensionless parameter groups reduce independent
parameters but do not eliminate empirical estimation needs.

\noindent\rule{\linewidth}{0.4pt}

% ============================================================
\section{Conclusion}\label{sec:conclusion}
% ============================================================

The H-Bar Model advances the case that depth, schema coherence,
attentional fidelity, executive control, metacognitive self-modelling,
and collective schema communication are formally independent variables
in agent training --- that optimising any subset without the others
produces systematically predictable failure modes --- and that the field
now has both the empirical tools and the formal framework to test this
claim rigorously.

\noindent\rule{\linewidth}{0.4pt}

% ============================================================
% Appendix A --- Notation Table
% ============================================================
\section*{Appendix A: Notation Reference}\label{app:notation}

\begin{center}
\resizebox{\textwidth}{!}{%
\small
\begin{tabular}{c l l l}
\toprule
\textbf{Symbol} & \textbf{Type} & \textbf{Domain} & \textbf{Description} \\
\midrule
$D$ & Set & Global & Knowledge domains, indexed by $d \in D$ \\
$M$ & Set & Global & Modalities $\{$text, visual, auditory, sensorimotor, symbolic$\}$ \\
$\Delta(d,t)$ & Function & $\R_+$ & Domain frontier (maximum achievable depth) \\
$\delta_A(d,t)$ & State var & $[0, \Delta(d,t)]$ & Parametric depth (structural complexity) \\
$\delta_A^{\text{relative}}(d,t)$ & Derived & $[0,1]$ & Relative depth (fraction of frontier achieved) \\
$\beta_A(d,t)$ & State var & $\R_+$ & Functional breadth (surface competence) \\
$\sigma_A(d,t)$ & State var & $[0,1]$ & Schema coherence (principled reorganisation degree) \\
$\tilde{\sigma}_A(d,t)$ & Estimate & $[0,1]$ & Stage 1 training-time proxy (multi-signal fusion) \\
$\hat{\sigma}_A(d,t)$ & Estimate & $[0,1]$ & Stage 2 evaluation-time ground-truth (OOD/ID ratio) \\
$\alpha_A(d,t)$ & State var & $[0,1]$ & Attentional fidelity (structure-directed attention) \\
$\hat{M}_A(d,t)$ & State var & $[0,1]$ & Self-model of schema coherence (metacognitive estimate) \\
$\zeta_A(d,t)$ & Derived & $[-1,1]$ & Calibration error ($\hat{M}_A - \sigma_A$) \\
$\Xi_A^P(t)$ & State var & $[0,1]$ & Executive planning quality \\
$\Xi_A^I(t)$ & State var & $[0,1]$ & Executive inhibition quality \\
$\Xi_A^F(t)$ & State var & $[0,1]$ & Executive cognitive flexibility \\
$\Psi_A(d_1,d_2,t)$ & Derived & $\R_+$ & Intersection activation (cross-domain discovery rate) \\
$q_A(d,t)$ & Derived & $[0,1]$ & Effective mastery quality ($\sigma_A \cdot \delta_A^{\text{relative}}$) \\
$M_A(t)$ & Set & $2^D$ & Mastery set (domains with high $\delta_A$ and $\sigma_A$) \\
$B_A(t)$ & Set & $2^D$ & Breadth set (domains with high $\beta_A$ but not mastered) \\
$D^*(d,t)$ & Derived & $2^d$ & Delegation gradient (skills safe for AI handoff) \\
$\lambda_c$ & Parameter & $\R_+$ & Parametric decay rate \\
$\lambda_c^{\text{eff}}$ & Derived & $\R_+$ & Effective decay (schema-mediated reduction) \\
$\lambda_f(d,t)$ & Parameter & $\R_+$ & Frontier obsolescence rate \\
$\eta(d,t)$ & Derived & $\R_+$ & Learning efficiency (Gompertz form) \\
$\eta_{\max}$ & Parameter & $\R_+$ & Maximum learning efficiency \\
$\rho$ & Parameter & $\R_+$ & Schema growth rate constant \\
$\epsilon_\sigma$ & Parameter & $\R_+$ & AI bypass suppression coefficient \\
$\Omega_{AI}(d,t)$ & Function & $\R_+$ & AI bypass risk (surface-reward pressure) \\
$P_A(d,t)$ & Function & $\R_+$ & Principled structure availability \\
$\gamma$ & Parameter & $\R_+$ & Attention growth rate constant \\
$\zeta_\alpha$ & Parameter & $\R_+$ & Surface-reward attention suppression \\
$R_A^{\text{surface}}(d,t)$ & Function & $\R_+$ & Surface-reward pressure \\
$\nu_M$ & Parameter & $\R_+$ & Self-model convergence rate \\
$\xi_M$ & Parameter & $\R_+$ & AI bypass metacognitive inflation \\
$\kappa_P, \kappa_I, \kappa_F$ & Parameters & $\R_+$ & Executive control convergence rates \\
$\mu_{AB}(d,t)$ & Derived & $[0,1]$ & Schema legibility (inter-agent) \\
$\Sigma_{A,B}(d_1,d_2,t)$ & Derived & $[0,1]$ & Collective schema field \\
$\tau_A(B,d,t)$ & Parameter & $[0,1]$ & Theory-of-mind coupling.  (Note: bare $\tau$ without subscript denotes the segmented-regression breakpoint time in Prediction~9.) \\
$\Theta_A(d,m_1,m_2,t)$ & Derived & $[0,1]$ & Cross-modal schema transfer \\
$\omega(m_1,m_2)$ & Parameter & $[0,1]$ & Modality similarity weight \\
$V_A(B,f,t)$ & Derived & $[0,1]$ & Benchmark validity function \\
$CI(B,f)$ & Metric & $[0,1]$ & Construct isolation \\
$FD(B)$ & Metric & $[0,1]$ & Format diversity \\
$DG(B)$ & Metric & $[0,1]$ & Difficulty gradient \\
$R_A(B,f,t)$ & Metric & $[0,1]$ & Benchmark reliability \\
$\sigma_{\text{critical}}$ & Derived & $[0,1]$ & Phase 1$\to$2 transition threshold \\
$R_0$ & Derived & $\R_+$ & Basic schema reproduction number \\
$\theta_I, \theta_\delta, \theta_\sigma, \theta_\beta$ & Parameters & $[0,1]$ & Threshold parameters \\
$\gamma_\sigma$ & Parameter & $[0,1]$ & Schema-mediated decay protection \\
$w_g, w_r, w_c$ & Parameters & $[0,1]$ & Signal fusion weights (Stage 1) \\
\bottomrule
\end{tabular}%
}
\end{center}

\textbf{Visual syntax conventions:}
\begin{itemize}[nosep]
\item No decoration (e.g., $\sigma_A$): true latent variable
\item Tilde (e.g., $\tilde{\sigma}_A$): training-time proxy estimate (Stage 1)
\item Hat (e.g., $\hat{\sigma}_A$): evaluation-time ground-truth estimate (Stage 2)
\item Superscript ``relative'' (e.g., $\delta_A^{\text{relative}}$): normalised to $[0,1]$
\item Superscript ``eff'' (e.g., $\lambda_c^{\text{eff}}$): effective/derived quantity
\item Subscript $A$: agent-relative quantity
\item Subscript $AI$: AI-system quantity
\end{itemize}

\noindent\rule{\linewidth}{0.4pt}

% ============================================================
% Bibliography
% ============================================================
\bibliographystyle{plainnat}
\bibliography{bibliography}

\end{document}
